\documentclass[a4paper,11pt,oneside,enabledeprecatedfontcommands]{scrbook}

% Pengaturan Font dan Bahasa untuk LuaLaTeX
\usepackage{fontspec}
\usepackage[indonesian]{babel}
\usepackage{amsmath,amssymb,amsfonts}
\usepackage{graphicx}
\usepackage[space]{grffile} % Mendukung nama file gambar yang mengandung spasi seperti "Tugas 7-002.png"
\usepackage{xcolor}
\usepackage{listings}
\usepackage{fancyvrb}
\usepackage{geometry}
\usepackage{hyperref}

\geometry{
    a4paper,
    top=2.5cm,
    bottom=2.5cm,
    left=2.5cm,
    right=2.5cm
}

% Definisi warna ala Euler Math Toolbox (EMT)
\definecolor{eulerpromptcolor}{rgb}{0.6,0.0,0.0}
\definecolor{euleroutputcolor}{rgb}{0.15,0.15,0.45}
\definecolor{eulerbg}{rgb}{0.97,0.97,0.98}
\definecolor{eulerborder}{rgb}{0.82,0.85,0.90}

% Pengaturan Listings untuk Kode EMT
\lstset{
    basicstyle=\ttfamily\small,
    breaklines=true,
    breakatwhitespace=false,
    columns=flexible,
    keepspaces=true,
    showstringspaces=false,
    frame=leftline,
    framerule=1.5pt,
    rulecolor=\color{eulerborder},
    backgroundcolor=\color{eulerbg},
    xleftmargin=10pt,
    aboveskip=6pt,
    belowskip=4pt
}

% Lingkungan khusus pengganti paket eumat agar kompatibel penuh dengan LuaLaTeX
\newenvironment{eulernotebook}{}{}

\newcommand{\eulerheading}[1]{\chapter{#1}}
\newcommand{\eulersubheading}[1]{\section{#1}}

\newenvironment{eulercomment}{\par\medskip\noindent}{\par\medskip}
\newenvironment{eulerformula}{}{}

\lstnewenvironment{eulerprompt}{
    \lstset{basicstyle=\ttfamily\small\color{eulerpromptcolor}}
}{}

\lstnewenvironment{euleroutput}{
    \lstset{
        basicstyle=\ttfamily\footnotesize\color{euleroutputcolor},
        backgroundcolor=\color{white},
        rulecolor=\color{lightgray}
    }
}{}

\lstnewenvironment{eulerudf}{
    \lstset{basicstyle=\ttfamily\small\color{teal!70!black}}
}{}

\lstnewenvironment{eulerttcomment}{
    \lstset{basicstyle=\ttfamily\small\color{darkgray}}
}{}

% Perintah penampil gambar yang memperbaiki spasi path dan mengabaikan tinggi 0
\newcommand{\eulerimg}[2]{%
    \ifnum#1>0
        \begin{center}
            \IfFileExists{#2}{%
                \includegraphics[width=0.72\textwidth,keepaspectratio]{{"#2"}}%
            }{%
                \fbox{\parbox{0.7\textwidth}{\centering\small\texttt{Gambar tidak ditemukan: #2}\\ \footnotesize Pastikan folder \texttt{images/} berada satu direktori dengan file \texttt{.tex}.}}%
            }%
        \end{center}
    \fi
}

\title{\textbf{Statistika dengan Euler Math Toolbox (EMT)}}
\author{Catatan Praktikum dan Latihan Statistika}
\date{\today}

\begin{document}
	
	\frontmatter
	\maketitle
	\tableofcontents
	
	\mainmatter
	\begin{eulernotebook}
		\eulerheading{EMT untuk Statistika}
		\begin{eulercomment}
			Pada notebook ini, kami mendemonstrasikan plot, uji, dan distribusi statistika utama di Euler.
			
			Mari kita mulai dengan beberapa statistika deskriptif. Ini bukan pengantar statistika. Jadi Anda
			mungkin memerlukan beberapa latar belakang untuk memahami detailnya.
			
			Misalkan terdapat pengukuran berikut. Kita ingin menghitung nilai rata-rata (mean) dan simpangan baku
			(standar deviasi) yang terukur.
		\end{eulercomment}
		\begin{eulerprompt}
			>M=[1000,1004,998,997,1002,1001,998,1004,998,997]; ...
			>median(M), mean(M), dev(M)
		\end{eulerprompt}
		\begin{euleroutput}
			999
			999.9
			2.72641400622
		\end{euleroutput}
		\begin{eulerprompt}
			>a&=max(M)-min(M); $a // range M
		\end{eulerprompt}
		\begin{eulerformula}
			\[
			0
			\]
		\end{eulerformula}
		\begin{eulercomment}
			Kita dapat membuat plot box-and-whiskers (boxplot) untuk data tersebut. Dalam kasus kita, tidak ada
			pencilan (outlier).
		\end{eulercomment}
		\begin{eulerprompt}
			>aspect(1.75); boxplot(M):
		\end{eulerprompt}
		\eulerimg{25}{images/Tugas 7-002.png}
		\begin{eulercomment}
			Kita menghitung probabilitas bahwa suatu nilai lebih besar dari 1005, dengan mengasumsikan
			nilai-nilai yang diukur berasal dari distribusi normal.
			
			Semua fungsi untuk distribusi di Euler diakhiri dengan \texttt{...dis} dan menghitung fungsi distribusi
			probabilitas kumulatif (CPF).
		\end{eulercomment}
		\begin{eulerformula}
			\[
			\text{normaldis}(x,m,d)=\int_{-\infty}^x \frac{1}{d\sqrt{2\pi}}e^{-\frac{1}{2}\left(\frac{t-m}{d}\right)^2}\,dt.
			\]
		\end{eulerformula}
		\begin{eulercomment}
			Kita mencetak hasilnya dalam \% dengan ketelitian 2 digit menggunakan fungsi \texttt{print}.
		\end{eulercomment}
		\begin{eulerprompt}
			>print((1-normaldis(1005,mean(M),dev(M)))*100,2,unit=" %")
		\end{eulerprompt}
		\begin{euleroutput}
			3.07 %
		\end{euleroutput}
		\begin{eulercomment}
			Untuk contoh berikutnya, kita mengasumsikan jumlah pria berikut dalam rentang ukuran tertentu.
		\end{eulercomment}
		\begin{eulerprompt}
			>r=155.5:4:187.5; v=[22,71,136,169,139,71,32,8];
		\end{eulerprompt}
		\begin{eulercomment}
			Berikut adalah plot dari distribusi tersebut.
		\end{eulercomment}
		\begin{eulerprompt}
			>plot2d(r,v,a=150,b=200,c=0,d=190,bar=1,style="\/"):
		\end{eulerprompt}
		\eulerimg{25}{images/Tugas 7-004.png}
		\begin{eulercomment}
			Kita dapat memasukkan data mentah tersebut ke dalam sebuah tabel.
			
			Tabel adalah metode untuk menyimpan data statistika. Tabel kita harus berisi tiga kolom: Awal
			rentang, akhir rentang, dan jumlah pria dalam rentang tersebut.
			
			Tabel dapat dicetak dengan judul kolom (header). Kita menggunakan vektor string untuk mengatur judul
			kolomnya.
		\end{eulercomment}
		\begin{eulerprompt}
			>T:=r[1:8]' | r[2:9]' | v'; writetable(T,labc=["BB","BA","Frek"])
		\end{eulerprompt}
		\begin{euleroutput}
			      BB        BA      Frek
			   155.5     159.5        22
			   159.5     163.5        71
			   163.5     167.5       136
			   167.5     171.5       169
			   171.5     175.5       139
			   175.5     179.5        71
			   179.5     183.5        32
			   183.5     187.5         8
		\end{euleroutput}
		\begin{eulercomment}
			Jika kita memerlukan nilai rata-rata dan statistik lainnya dari ukuran tersebut, kita perlu
			menghitung titik tengah dari rentang-rentang tersebut. Kita dapat menggunakan dua kolom pertama dari
			tabel kita untuk ini.
			
			Simbol \texttt{|} digunakan untuk memisahkan kolom, fungsi \texttt{writetable} digunakan untuk menulis tabel,
			dengan opsi \texttt{labc} untuk menentukan judul kolom.
		\end{eulercomment}
		\begin{eulerprompt}
			>(T[,1]+T[,2])/2 // titik tengah setiap interval
		\end{eulerprompt}
		\begin{euleroutput}
			157.5 
			161.5 
			165.5 
			169.5 
			173.5 
			177.5 
			181.5 
			185.5 
		\end{euleroutput}
		\begin{eulercomment}
			Namun, lebih mudah untuk melipat (fold) rentang tersebut dengan vektor \texttt{[1/2,1/2]}.
		\end{eulercomment}
		\begin{eulerprompt}
			>M=fold(r,[0.5,0.5])
		\end{eulerprompt}
		\begin{euleroutput}
			[157.5,  161.5,  165.5,  169.5,  173.5,  177.5,  181.5,  185.5]
		\end{euleroutput}
		\begin{eulercomment}
			Sekarang kita dapat menghitung rata-rata dan simpangan dari sampel dengan frekuensi yang diberikan.
		\end{eulercomment}
		\begin{eulerprompt}
			>{m,d}=meandev(M,v); m, d,
		\end{eulerprompt}
		\begin{euleroutput}
			169.901234568
			5.98912964449
		\end{euleroutput}
		\begin{eulercomment}
			Mari kita tambahkan distribusi normal dari nilai-nilai tersebut ke plot batang di atas. Rumus
			distribusi normal dengan rata-rata $m$ dan simpangan baku $d$ adalah:
		\end{eulercomment}
		\begin{eulerformula}
			\[
			y=\frac{1}{d\sqrt{2\pi}}e^{\frac{-(x-m)^2}{2d^2}}.
			\]
		\end{eulerformula}
		\begin{eulercomment}
			Karena nilainya berada di antara 0 dan 1, untuk memplotnya pada plot batang, nilainya harus dikalikan
			dengan 4 kali jumlah total data.
		\end{eulercomment}
		\begin{eulerprompt}
			>plot2d("qnormal(x,m,d)*sum(v)*4", ...
			>  xmin=min(r),xmax=max(r),thickness=3,add=1):
		\end{eulerprompt}
		\eulerimg{25}{images/Tugas 7-006.png}
		
		\eulerheading{Tabel}
		\begin{eulercomment}
			Di dalam direktori notebook ini Anda akan menemukan sebuah file berisi tabel. Data tersebut
			merepresentasikan hasil survei. Berikut adalah empat baris pertama dari file tersebut. Data ini
			berasal dari buku daring berbahasa Jerman ``Einführung in die Statistik mit R'' karya A. Handl.
		\end{eulercomment}
		\begin{eulerprompt}
			>printfile("table.dat",4);
		\end{eulerprompt}
		\begin{euleroutput}
			Person Sex Age Titanic Evaluation Tip Problem
			1 m 30 n . 1.80 n
			2 f 23 y g 1.80 n
			3 f 26 y g 1.80 y
		\end{euleroutput}
		\begin{eulercomment}
			Tabel ini berisi 7 kolom angka atau token (string). Kita ingin membaca tabel dari file tersebut.
			Pertama, kita menggunakan terjemahan kita sendiri untuk token-token tersebut.
			
			Untuk itu, kita mendefinisikan himpunan token. Fungsi \texttt{strtokens()} mengambil vektor string token dari
			string yang diberikan.
		\end{eulercomment}
		\begin{eulerprompt}
			>mf:=["m","f"]; yn:=["y","n"]; ev:=strtokens("g vg m b vb");
		\end{eulerprompt}
		\begin{eulercomment}
			Sekarang kita membaca tabel dengan terjemahan ini.
			
			Argumen \texttt{tok2}, \texttt{tok4}, dll. adalah terjemahan dari kolom-kolom tabel. Argumen-argumen ini tidak ada
			dalam daftar parameter \texttt{readtable()}, sehingga Anda perlu memberikannya dengan \texttt{:=}.
		\end{eulercomment}
		\begin{eulerprompt}
			>{MT,hd}=readtable("table.dat",tok2:=mf,tok4:=yn,tok5:=ev,tok7:=yn);
			>load over statistics;
		\end{eulerprompt}
		\begin{eulercomment}
			Untuk pencetakan, kita perlu menentukan himpunan token yang sama. Kita hanya mencetak sepuluh baris
			pertama.
		\end{eulercomment}
		\begin{eulerprompt}
			>writetable(MT[1:10],labc=hd,wc=5,tok2:=mf,tok4:=yn,tok5:=ev,tok7:=yn);
		\end{eulerprompt}
		\begin{euleroutput}
			 Person  Sex  Age Titanic Evaluation  Tip Problem
			      1    m   30       n          .  1.8       n
			      2    f   23       y          g  1.8       n
			      3    f   26       y          g  1.8       y
			      4    m   33       n          .  2.8       n
			      5    m   37       n          .  1.8       n
			      6    m   28       y          g  2.8       y
			      7    f   31       y         vg  2.8       n
			      8    m   23       n          .  0.8       n
			      9    f   24       y         vg  1.8       y
			     10    m   26       n          .  1.8       n
		\end{euleroutput}
		\begin{eulercomment}
			Tanda titik \texttt{"."} merepresentasikan nilai yang tidak tersedia.
			
			Jika kita tidak ingin menentukan token untuk terjemahan terlebih dahulu, kita hanya perlu menentukan
			kolom mana yang berisi token dan bukan angka.
		\end{eulercomment}
		\begin{eulerprompt}
			>ctok=[2,4,5,7]; {MT,hd,tok}=readtable("table.dat",ctok=ctok);
		\end{eulerprompt}
		\begin{eulercomment}
			Fungsi \texttt{readtable()} sekarang mengembalikan himpunan token.
		\end{eulercomment}
		\begin{eulerprompt}
			>tok
		\end{eulerprompt}
		\begin{euleroutput}
			m
			n
			f
			y
			g
			vg
		\end{euleroutput}
		\begin{eulercomment}
			Tabel tersebut berisi entri dari file dengan token yang telah diterjemahkan menjadi angka.
			
			String khusus \texttt{NA="."} ditafsirkan sebagai ``Not Available'' (Tidak Tersedia), dan mendapatkan nilai \texttt{NAN}
			(not a number) dalam tabel. Terjemahan ini dapat diubah dengan parameter \texttt{NA} dan \texttt{NAval}.
		\end{eulercomment}
		\begin{eulerprompt}
			>MT[1]
		\end{eulerprompt}
		\begin{euleroutput}
			[1,  1,  30,  2,  NAN,  1.8,  2]
		\end{euleroutput}
		\begin{eulercomment}
			Berikut adalah isi tabel dengan angka yang belum diterjemahkan.
		\end{eulercomment}
		\begin{eulerprompt}
			>writetable(MT,wc=5)
		\end{eulerprompt}
		\begin{euleroutput}
			    1    1   30    2    .  1.8    2
			    2    3   23    4    5  1.8    2
			    3    3   26    4    5  1.8    4
			    4    1   33    2    .  2.8    2
			    5    1   37    2    .  1.8    2
			    6    1   28    4    5  2.8    4
			    7    3   31    4    6  2.8    2
			    8    1   23    2    .  0.8    2
			    9    3   24    4    6  1.8    4
			   10    1   26    2    .  1.8    2
			   11    3   23    4    6  1.8    4
			   12    1   32    4    5  1.8    2
			   13    1   29    4    6  1.8    4
			   14    3   25    4    5  1.8    4
			   15    3   31    4    5  0.8    2
			   16    1   26    4    5  2.8    2
			   17    1   37    2    .  3.8    2
			   18    1   38    4    5    .    2
			   19    3   29    2    .  3.8    2
			   20    3   28    4    6  1.8    2
			   21    3   28    4    1  2.8    4
			   22    3   28    4    6  1.8    4
			   23    3   38    4    5  2.8    2
			   24    3   27    4    1  1.8    4
			   25    1   27    2    .  2.8    4
		\end{euleroutput}
		\begin{eulercomment}
			Untuk kemudahan, Anda dapat memasukkan keluaran (output) dari \texttt{readtable()} ke dalam sebuah daftar
			(list).
		\end{eulercomment}
		\begin{eulerprompt}
			>Table={{readtable("table.dat",ctok=ctok)}};
		\end{eulerprompt}
		\begin{eulercomment}
			Dengan menggunakan kolom token yang sama dan token yang dibaca dari file, kita dapat mencetak tabel
			tersebut. Kita dapat menentukan \texttt{ctok}, \texttt{tok}, dll. atau menggunakan daftar \texttt{Table}.
		\end{eulercomment}
		\begin{eulerprompt}
			>writetable(Table,ctok=ctok,wc=5);
		\end{eulerprompt}
		\begin{euleroutput}
			 Person  Sex  Age Titanic Evaluation  Tip Problem
			      1    m   30       n          .  1.8       n
			      2    f   23       y          g  1.8       n
			      3    f   26       y          g  1.8       y
			      4    m   33       n          .  2.8       n
			      5    m   37       n          .  1.8       n
			      6    m   28       y          g  2.8       y
			      7    f   31       y         vg  2.8       n
			      8    m   23       n          .  0.8       n
			      9    f   24       y         vg  1.8       y
			     10    m   26       n          .  1.8       n
			     11    f   23       y         vg  1.8       y
			     12    m   32       y          g  1.8       n
			     13    m   29       y         vg  1.8       y
			     14    f   25       y          g  1.8       y
			     15    f   31       y          g  0.8       n
			     16    m   26       y          g  2.8       n
			     17    m   37       n          .  3.8       n
			     18    m   38       y          g    .       n
			     19    f   29       n          .  3.8       n
			     20    f   28       y         vg  1.8       n
			     21    f   28       y          m  2.8       y
			     22    f   28       y         vg  1.8       y
			     23    f   38       y          g  2.8       n
			     24    f   27       y          m  1.8       y
			     25    m   27       n          .  2.8       y
		\end{euleroutput}
		\begin{eulercomment}
			Fungsi \texttt{tablecol()} mengembalikan nilai kolom dari tabel, dengan melewati baris mana pun yang memiliki
			nilai \texttt{NAN} (\texttt{"."} dalam file), serta indeks kolom yang berisi nilai-nilai tersebut.
		\end{eulercomment}
		\begin{eulerprompt}
			>{c,i}=tablecol(MT,[5,6]);
		\end{eulerprompt}
		\begin{eulercomment}
			Kita dapat menggunakan ini untuk mengekstrak kolom dari tabel untuk dijadikan tabel baru.
		\end{eulercomment}
		\begin{eulerprompt}
			>j=[1,5,6]; writetable(MT[i,j],labc=hd[j],ctok=[2],tok=tok)
		\end{eulerprompt}
		\begin{euleroutput}
			    Person Evaluation       Tip
			         2          g       1.8
			         3          g       1.8
			         6          g       2.8
			         7         vg       2.8
			         9         vg       1.8
			        11         vg       1.8
			        12          g       1.8
			        13         vg       1.8
			        14          g       1.8
			        15          g       0.8
			        16          g       2.8
			        20         vg       1.8
			        21          m       2.8
			        22         vg       1.8
			        23          g       2.8
			        24          m       1.8
		\end{euleroutput}
		\begin{eulercomment}
			Tentu saja, kita perlu mengekstrak tabel itu sendiri dari daftar \texttt{Table} dalam kasus ini.
		\end{eulercomment}
		\begin{eulerprompt}
			>MT=Table[1];
		\end{eulerprompt}
		\begin{eulercomment}
			Tentu saja, kita juga dapat menggunakannya untuk menentukan nilai rata-rata dari suatu kolom atau
			nilai statistika lainnya.
		\end{eulercomment}
		\begin{eulerprompt}
			>mean(tablecol(MT,6))
		\end{eulerprompt}
		\begin{euleroutput}
			2.175
		\end{euleroutput}
		\begin{eulercomment}
			Fungsi \texttt{getstatistics()} mengembalikan elemen-elemen dalam vektor beserta jumlah kemunculannya. Kita
			menerapkannya pada nilai \texttt{"m"} dan \texttt{"f"} di kolom kedua tabel kita.
		\end{eulercomment}
		\begin{eulerprompt}
			>{xu,count}=getstatistics(tablecol(MT,2)); xu, count,
		\end{eulerprompt}
		\begin{euleroutput}
			[1,  3]
			[12,  13]
		\end{euleroutput}
		\begin{eulercomment}
			Kita dapat mencetak hasilnya dalam tabel baru.
		\end{eulercomment}
		\begin{eulerprompt}
			>writetable(count',labr=tok[xu])
		\end{eulerprompt}
		\begin{euleroutput}
			         m        12
			         f        13
		\end{euleroutput}
		\begin{eulercomment}
			Fungsi \texttt{selecttable()} mengembalikan tabel baru dengan nilai-nilai dalam satu kolom yang dipilih dari
			vektor indeks. Pertama, kita mencari indeks dari dua nilai kita dalam tabel token.
		\end{eulercomment}
		\begin{eulerprompt}
			>v:=indexof(tok,["g","vg"])
		\end{eulerprompt}
		\begin{euleroutput}
			[5,  6]
		\end{euleroutput}
		\begin{eulercomment}
			Sekarang kita dapat memilih baris-baris tabel yang memiliki salah satu nilai dalam \texttt{v} pada kolom
			ke-5-nya.
		\end{eulercomment}
		\begin{eulerprompt}
			>MT1:=MT[selectrows(MT,5,v)]; i:=sortedrows(MT1,5);
		\end{eulerprompt}
		\begin{eulercomment}
			Sekarang kita dapat mencetak tabel tersebut, dengan nilai yang telah diekstrak dan diurutkan pada
			kolom ke-5.
		\end{eulercomment}
		\begin{eulerprompt}
			>writetable(MT1[i],labc=hd,ctok=ctok,tok=tok,wc=7);
		\end{eulerprompt}
		\begin{euleroutput}
			 Person    Sex    Age Titanic Evaluation    Tip Problem
			      2      f     23       y          g    1.8       n
			      3      f     26       y          g    1.8       y
			      6      m     28       y          g    2.8       y
			     18      m     38       y          g      .       n
			     16      m     26       y          g    2.8       n
			     15      f     31       y          g    0.8       n
			     12      m     32       y          g    1.8       n
			     23      f     38       y          g    2.8       n
			     14      f     25       y          g    1.8       y
			      9      f     24       y         vg    1.8       y
			      7      f     31       y         vg    2.8       n
			     20      f     28       y         vg    1.8       n
			     22      f     28       y         vg    1.8       y
			     13      m     29       y         vg    1.8       y
			     11      f     23       y         vg    1.8       y
		\end{euleroutput}
		\begin{eulercomment}
			Untuk statistik berikutnya, kita ingin menghubungkan dua kolom dari tabel. Jadi kita mengekstrak
			kolom 2 dan 4 lalu mengurutkan tabelnya.
		\end{eulercomment}
		\begin{eulerprompt}
			>i=sortedrows(MT,[2,4]);  ...
			>  writetable(tablecol(MT[i],[2,4])',ctok=[1,2],tok=tok)
		\end{eulerprompt}
		\begin{euleroutput}
			         m         n
			         m         n
			         m         n
			         m         n
			         m         n
			         m         n
			         m         n
			         m         y
			         m         y
			         m         y
			         m         y
			         m         y
			         f         n
			         f         y
			         f         y
			         f         y
			         f         y
			         f         y
			         f         y
			         f         y
			         f         y
			         f         y
			         f         y
			         f         y
			         f         y
		\end{euleroutput}
		\begin{eulercomment}
			Dengan \texttt{getstatistics()}, kita juga dapat menghubungkan jumlah kemunculan di dua kolom tabel satu sama
			lain.
		\end{eulercomment}
		\begin{eulerprompt}
			>MT24=tablecol(MT,[2,4]); ...
			>{xu1,xu2,count}=getstatistics(MT24[1],MT24[2]); ...
			>writetable(count,labr=tok[xu1],labc=tok[xu2])
		\end{eulerprompt}
		\begin{euleroutput}
			                   n         y
			         m         7         5
			         f         1        12
		\end{euleroutput}
		\begin{eulercomment}
			Sebuah tabel dapat ditulis ke dalam file.
		\end{eulercomment}
		\begin{eulerprompt}
			>filename="test.dat"; ...
			>writetable(count,labr=tok[xu1],labc=tok[xu2],file=filename);
		\end{eulerprompt}
		\begin{eulercomment}
			Kemudian kita dapat membaca tabel tersebut dari file.
		\end{eulercomment}
		\begin{eulerprompt}
			>{MT2,hd,tok2,hdr}=readtable(filename,>clabs,>rlabs); ...
			>writetable(MT2,labr=hdr,labc=hd)
		\end{eulerprompt}
		\begin{euleroutput}
			                   n         y
			         m         7         5
			         f         1        12
		\end{euleroutput}
		\begin{eulercomment}
			Dan menghapus file tersebut.
		\end{eulercomment}
		\begin{eulerprompt}
			>fileremove(filename);
		\end{eulerprompt}
		
		\eulerheading{Distribusi}
		\begin{eulercomment}
			Dengan \texttt{plot2d}, terdapat metode yang sangat mudah untuk memplot distribusi data eksperimen.
		\end{eulercomment}
		\begin{eulerprompt}
			>p=normal(1,1000); // 1000 sampel acak berdistribusi normal p
			>plot2d(p,distribution=20,style="\/"); // plot sampel acak p
			>plot2d("qnormal(x,0,1)",add=1): // tambahkan plot distribusi normal standar
		\end{eulerprompt}
		\eulerimg{25}{images/Tugas 7-007.png}
		\begin{eulercomment}
			Harap perhatikan perbedaan antara plot batang (sampel) dan kurva normal (distribusi sebenarnya).
			Masukkan kembali ketiga perintah tersebut untuk melihat hasil pengambilan sampel lainnya.
			
			Berikut adalah perbandingan 10 simulasi dari 1000 nilai berdistribusi normal menggunakan apa yang
			disebut plot kotak (box plot). Plot ini menunjukkan median, kuartil 25\% dan 75\%, nilai minimum dan
			maksimum, serta pencilan (outlier).
		\end{eulercomment}
		\begin{eulerprompt}
			>p=normal(10,1000); boxplot(p):
		\end{eulerprompt}
		\eulerimg{25}{images/Tugas 7-008.png}
		\begin{eulercomment}
			Untuk menghasilkan bilangan bulat acak, Euler memiliki \texttt{intrandom}. Mari kita simulasikan lemparan dadu
			dan plot distribusinya.
			
			Kita menggunakan fungsi \texttt{getmultiplicities(v,x)}, yang menghitung seberapa sering elemen-elemen \texttt{v}
			muncul dalam \texttt{x}. Kemudian kita memplot hasilnya menggunakan \texttt{columnsplot()}.
		\end{eulercomment}
		\begin{eulerprompt}
			>k=intrandom(1,6000,6);  ...
			>columnsplot(getmultiplicities(1:6,k));  ...
			>ygrid(1000,color=red):
		\end{eulerprompt}
		\eulerimg{25}{images/Tugas 7-009.png}
		\begin{eulercomment}
			Meskipun \texttt{intrandom(n,m,k)} mengembalikan bilangan bulat berdistribusi seragam dari 1 hingga $k$,
			dimungkinkan untuk menggunakan distribusi bilangan bulat lainnya dengan \texttt{randpint()}.
			
			Pada contoh berikut, probabilitas untuk 1, 2, 3 masing-masing adalah 0.4, 0.1, 0.5.
		\end{eulercomment}
		\begin{eulerprompt}
			>randpint(1,1000,[0.4,0.1,0.5]); getmultiplicities(1:3,%)
		\end{eulerprompt}
		\begin{euleroutput}
			[378,  102,  520]
		\end{euleroutput}
		\begin{eulercomment}
			Euler dapat menghasilkan nilai acak dari lebih banyak distribusi. Silakan lihat pada referensi.
			
			Misalnya, kita mencoba distribusi eksponensial. Variabel acak kontinu $X$ dikatakan memiliki distribusi
			eksponensial, jika PDF-nya diberikan oleh
		\end{eulercomment}
		\begin{eulerformula}
			\[
			f_X(x)=\lambda e^{-\lambda x},\quad x>0,\quad \lambda>0,
			\]
		\end{eulerformula}
		\begin{eulercomment}
			dengan parameter
		\end{eulercomment}
		\begin{eulerformula}
			\[
			\lambda=\frac{1}{\mu},\quad \mu \text{ adalah nilai rata-rata, dan dinotasikan dengan } X \sim \text{Exponential}(\lambda).
			\]
		\end{eulerformula}
		\begin{eulerprompt}
			>plot2d(randexponential(1,1000,2),>distribution):
		\end{eulerprompt}
		\eulerimg{25}{images/Tugas 7-012.png}
		\begin{eulercomment}
			Untuk banyak distribusi, Euler dapat menghitung fungsi distribusi dan inversnya.
		\end{eulercomment}
		\begin{eulerprompt}
			>plot2d("normaldis",-4,4): 
		\end{eulerprompt}
		\eulerimg{25}{images/Tugas 7-013.png}
		\begin{eulercomment}
			Berikut adalah salah satu cara untuk memplot kuantil.
		\end{eulercomment}
		\begin{eulerprompt}
			>plot2d("qnormal(x,1,1.5)",-4,6);  ...
			>plot2d("qnormal(x,1,1.5)",a=2,b=5,>add,>filled):
		\end{eulerprompt}
		\eulerimg{25}{images/Tugas 7-014.png}
		\begin{eulerformula}
			\[
			\text{normaldis}(x,m,d)=\int_{-\infty}^x \frac{1}{d\sqrt{2\pi}}e^{-\frac{1}{2}\left(\frac{t-m}{d}\right)^2}\,dt.
			\]
		\end{eulerformula}
		\begin{eulercomment}
			Probabilitas berada di area hijau adalah sebagai berikut.
		\end{eulercomment}
		\begin{eulerprompt}
			>normaldis(5,1,1.5)-normaldis(2,1,1.5)
		\end{eulerprompt}
		\begin{euleroutput}
			0.248662156979
		\end{euleroutput}
		\begin{eulercomment}
			Ini dapat dihitung secara numerik dengan integral berikut.
		\end{eulercomment}
		\begin{eulerformula}
			\[
			\int_2^5 \frac{1}{1.5\sqrt{2\pi}}e^{-\frac{1}{2}\left(\frac{x-1}{1.5}\right)^2}\,dx.
			\]
		\end{eulerformula}
		\begin{eulerprompt}
			>gauss("qnormal(x,1,1.5)",2,5)
		\end{eulerprompt}
		\begin{euleroutput}
			0.248662156979
		\end{euleroutput}
		\begin{eulercomment}
			Mari kita bandingkan distribusi binomial dengan distribusi normal yang memiliki rata-rata dan
			simpangan yang sama. Fungsi \texttt{invbindis()} menyelesaikan interpolasi linear di antara nilai-nilai
			bilangan bulat.
		\end{eulercomment}
		\begin{eulerprompt}
			>invbindis(0.95,1000,0.5), invnormaldis(0.95,500,0.5*sqrt(1000))
		\end{eulerprompt}
		\begin{euleroutput}
			525.516721219
			526.007419394
		\end{euleroutput}
		\begin{eulercomment}
			Fungsi \texttt{qchidis()} adalah kepadatan (densitas) dari distribusi chi-kuadrat. Seperti biasa, Euler memetakan
			vektor ke fungsi ini. Dengan demikian kita mendapatkan plot semua distribusi chi-kuadrat dengan
			derajat 5 hingga 50 dengan mudah melalui cara berikut.
		\end{eulercomment}
		\begin{eulerprompt}
			>plot2d("qchidis(x,(5:5:50)')",0,50):
		\end{eulerprompt}
		\eulerimg{25}{images/Tugas 7-017.png}
		\begin{eulercomment}
			Euler memiliki fungsi-fungsi yang akurat untuk mengevaluasi distribusi. Mari kita periksa \texttt{chidis()}
			dengan sebuah integral.
			
			Penamaannya diupayakan agar konsisten. Misalnya:
			\begin{itemize}
				\item distribusi chi-kuadrat adalah \texttt{chidis()},
				\item fungsi inversnya adalah \texttt{invchidis()},
				\item kepadatannya adalah \texttt{qchidis()}.
			\end{itemize}
			
			Komplemen dari distribusi tersebut (ekor atas) adalah \texttt{chicdis()}.
		\end{eulercomment}
		\begin{eulerprompt}
			>chidis(1.5,2), integrate("qchidis(x,2)",0,1.5)
		\end{eulerprompt}
		\begin{euleroutput}
			0.527633447259
			0.527633447259
		\end{euleroutput}
		
		\eulerheading{Distribusi Diskrit}
		\begin{eulercomment}
			Untuk mendefinisikan distribusi diskrit Anda sendiri, Anda dapat menggunakan metode berikut.
			
			Pertama, kita mengatur fungsi distribusinya.
		\end{eulercomment}
		\begin{eulerprompt}
			>wd = 0|((1:6)+[-0.01,0.01,0,0,0,0])/6
		\end{eulerprompt}
		\begin{euleroutput}
			[0,  0.165,  0.335,  0.5,  0.666667,  0.833333,  1]
		\end{euleroutput}
		\begin{eulercomment}
			Artinya adalah bahwa dengan probabilitas \texttt{wd[i+1]-wd[i]} kita menghasilkan nilai acak $i$.
			
			Ini hampir merupakan distribusi seragam. Mari kita definisikan pembangkit bilangan acak untuk ini.
			Fungsi \texttt{find(v,x)} mencari nilai $x$ dalam vektor $v$. Fungsi ini juga bekerja untuk vektor $x$.
		\end{eulercomment}
		\begin{eulerprompt}
			>function wrongdice (n,m) := find(wd,random(n,m))
		\end{eulerprompt}
		\begin{eulercomment}
			Kesalahannya sangat halus sehingga kita baru melihatnya dengan iterasi yang sangat banyak.
		\end{eulercomment}
		\begin{eulerprompt}
			>columnsplot(getmultiplicities(1:6,wrongdice(1,1000000))):
		\end{eulerprompt}
		\eulerimg{25}{images/Tugas 7-018.png}
		\begin{eulercomment}
			Berikut adalah fungsi sederhana untuk memeriksa distribusi seragam dari nilai $1\dots K$ dalam $v$. Kita
			menerima hasilnya jika untuk semua frekuensi berlaku
		\end{eulercomment}
		\begin{eulerformula}
			\[
			\left|f_i-\frac{1}{K}\right| < \frac{\delta}{\sqrt{n}}.
			\]
		\end{eulerformula}
		\begin{eulerprompt}
			>function checkrandom (v, delta=1) ...
		\end{eulerprompt}
		\begin{eulerudf}
			K=max(v); n=cols(v);
			fr=getfrequencies(v,1:K);
			return max(fr/n-1/K)<delta/sqrt(n);
			endfunction
		\end{eulerudf}
		\begin{eulercomment}
			Memang, fungsi tersebut menolak distribusi seragam.
		\end{eulercomment}
		\begin{eulerprompt}
			>checkrandom(wrongdice(1,1000000))
		\end{eulerprompt}
		\begin{euleroutput}
			0
		\end{euleroutput}
		\begin{eulercomment}
			Dan fungsi tersebut menerima pembangkit acak bawaan.
		\end{eulercomment}
		\begin{eulerprompt}
			>checkrandom(intrandom(1,1000000,6))
		\end{eulerprompt}
		\begin{euleroutput}
			1
		\end{euleroutput}
		\begin{eulercomment}
			Kita dapat menghitung distribusi binomial. Pertama ada \texttt{bindis()}, yang mengembalikan probabilitas
			$i$ atau kurang keberhasilan dari $n$ percobaan.
		\end{eulercomment}
		\begin{eulerprompt}
			>bindis(410,1000,0.4)
		\end{eulerprompt}
		\begin{euleroutput}
			0.751401349654
		\end{euleroutput}
		\begin{eulercomment}
			Fungsi invers Beta digunakan untuk menghitung interval kepercayaan Clopper-Pearson untuk parameter $p$.
			Tingkat default-nya adalah $\alpha$.
			
			Arti dari interval ini adalah bahwa jika $p$ berada di luar interval tersebut, hasil yang diamati yaitu
			410 dari 1000 adalah jarang terjadi.
		\end{eulercomment}
		\begin{eulerprompt}
			>clopperpearson(410,1000)
		\end{eulerprompt}
		\begin{euleroutput}
			[0.37932,  0.441212]
		\end{euleroutput}
		\begin{eulercomment}
			Perintah-perintah berikut adalah cara langsung untuk mendapatkan hasil di atas. Namun untuk $n$ yang
			besar, penjumlahan langsung tidak akurat dan lambat.
		\end{eulercomment}
		\begin{eulerprompt}
			>p=0.4; i=0:410; n=1000; sum(bin(n,i)*p^i*(1-p)^(n-i))
		\end{eulerprompt}
		\begin{euleroutput}
			0.751401349655
		\end{euleroutput}
		\begin{eulercomment}
			Omong-omong, \texttt{invbindis()} menghitung invers dari distribusi kumulatif binomial.
		\end{eulercomment}
		\begin{eulerprompt}
			>invbindis(0.75,1000,0.4)
		\end{eulerprompt}
		\begin{euleroutput}
			409.932733047
		\end{euleroutput}
		\begin{eulercomment}
			Dalam permainan Bridge, kita mengasumsikan 5 kartu yang tersisa (dari 52) di dua tangan (26 kartu).
			Mari kita hitung probabilitas distribusi yang lebih buruk dari 3:2 (misalnya 0:5, 1:4, 4:1 atau 5:0).
		\end{eulercomment}
		\begin{eulerprompt}
			>2*hypergeomsum(1,5,13,26)
		\end{eulerprompt}
		\begin{euleroutput}
			0.321739130435
		\end{euleroutput}
		\begin{eulercomment}
			Terdapat juga simulasi distribusi multinomial.
		\end{eulercomment}
		\begin{eulerprompt}
			>randmultinomial(10,1000,[0.4,0.1,0.5])
		\end{eulerprompt}
		\begin{euleroutput}
			      381           100           519 
			      376            91           533 
			      417            80           503 
			      440            94           466 
			      406           112           482 
			      408            94           498 
			      395           107           498 
			      399            96           505 
			      428            87           485 
			      400            99           501 
		\end{euleroutput}
		
		\eulerheading{Memplot Data}
		\begin{eulercomment}
			Untuk memplot data, kita mencoba hasil pemilu Jerman sejak tahun 1990, yang diukur dalam jumlah
			kursi.
		\end{eulercomment}
		\begin{eulerprompt}
			>BW := [ ...
			>1990,662,319,239,79,8,17; ...
			>1994,672,294,252,47,49,30; ...
			>1998,669,245,298,43,47,36; ...
			>2002,603,248,251,47,55,2; ...
			>2005,614,226,222,61,51,54; ...
			>2009,622,239,146,93,68,76; ...
			>2013,631,311,193,0,63,64];
		\end{eulerprompt}
		\begin{eulercomment}
			Untuk partai-partainya, kita menggunakan string nama.
		\end{eulercomment}
		\begin{eulerprompt}
			>P:=["CDU/CSU","SPD","FDP","Gr","Li"];
		\end{eulerprompt}
		\begin{eulercomment}
			Mari kita cetak persentasenya dengan rapi.
			
			Pertama kita mengekstrak kolom-kolom yang diperlukan. Kolom 3 hingga 7 adalah kursi masing-masing
			partai, dan kolom 2 adalah jumlah total kursi. Kolom 1 adalah tahun pemilu.
		\end{eulercomment}
		\begin{eulerprompt}
			>BT:=BW[,3:7]; BT:=BT/sum(BT); YT:=BW[,1]';
		\end{eulerprompt}
		\begin{eulercomment}
			Kemudian kita mencetak statistik dalam bentuk tabel. Kita menggunakan nama partai sebagai judul
			kolom, dan tahun sebagai judul baris. Lebar default untuk kolom adalah \texttt{wc=10}, tetapi kita lebih
			menyukai tampilan yang lebih rapat. Kolom akan diperlebar untuk label kolom jika diperlukan.
		\end{eulercomment}
		\begin{eulerprompt}
			>writetable(BT*100,wc=6,dc=0,>fixed,labc=P,labr=YT)
		\end{eulerprompt}
		\begin{euleroutput}
			       CDU/CSU   SPD   FDP    Gr    Li
			  1990      48    36    12     1     3
			  1994      44    38     7     7     4
			  1998      37    45     6     7     5
			  2002      41    42     8     9     0
			  2005      37    36    10     8     9
			  2009      38    23    15    11    12
			  2013      49    31     0    10    10
		\end{euleroutput}
		\begin{eulercomment}
			Perkalian matriks berikut mengekstrak jumlah persentase dari dua partai besar yang menunjukkan bahwa
			partai-partai kecil telah memperoleh pijakan di parlemen hingga tahun 2009.
		\end{eulercomment}
		\begin{eulerprompt}
			>BT1:=(BT.[1;1;0;0;0])'*100
		\end{eulerprompt}
		\begin{euleroutput}
			[84.29,  81.25,  81.1659,  82.7529,  72.9642,  61.8971,  79.8732]
		\end{euleroutput}
		\begin{eulercomment}
			Terdapat juga plot statistika sederhana. Kita menggunakannya untuk menampilkan garis dan titik secara
			bersamaan. Alternatifnya adalah memanggil \texttt{plot2d} dua kali dengan \texttt{>add}.
		\end{eulercomment}
		\begin{eulerprompt}
			>statplot(YT,BT1,"b"):
		\end{eulerprompt}
		\eulerimg{25}{images/Tugas 7-020.png}
		\begin{eulercomment}
			Mendefinisikan beberapa warna untuk masing-masing partai.
		\end{eulercomment}
		\begin{eulerprompt}
			>CP:=[rgb(0.5,0.5,0.5),red,yellow,green,rgb(0.8,0,0)];
		\end{eulerprompt}
		\begin{eulercomment}
			Sekarang kita dapat memplot hasil pemilu 2009 dan perubahannya ke dalam satu plot menggunakan \texttt{figure}.
			Kita dapat menambahkan vektor kolom ke setiap plot.
		\end{eulercomment}
		\begin{eulerprompt}
			>figure(2,1);  ...
			>figure(1); columnsplot(BW[6,3:7],P,color=CP); ...
			>figure(2); columnsplot(BW[6,3:7]-BW[5,3:7],P,color=CP);  ...
			>figure(0):
		\end{eulerprompt}
		\eulerimg{25}{images/Tugas 7-021.png}
		\begin{eulercomment}
			Plot data menggabungkan baris-baris data statistika dalam satu plot.
		\end{eulercomment}
		\begin{eulerprompt}
			>J:=BW[,1]'; DP:=BW[,3:7]'; ...
			>dataplot(YT,BT',color=CP);  ...
			>labelbox(P,colors=CP,styles="[]",>points,w=0.2,x=0.3,y=0.4):
		\end{eulerprompt}
		\eulerimg{25}{images/Tugas 7-022.png}
		\begin{eulercomment}
			Plot kolom 3D menunjukkan baris-baris data statistika dalam bentuk kolom. Kita memberikan label untuk
			baris dan kolom. \texttt{angle} adalah sudut pandang.
		\end{eulercomment}
		\begin{eulerprompt}
			>columnsplot3d(BT,scols=P,srows=YT, ...
			>  angle=30,ccols=CP):
		\end{eulerprompt}
		\eulerimg{25}{images/Tugas 7-023.png}
		\begin{eulercomment}
			Representasi lainnya adalah plot mosaik. Perhatikan bahwa kolom-kolom plot di sini merepresentasikan
			kolom-kolom matriks. Karena panjangnya label CDU/CSU, kita mengambil jendela yang lebih kecil dari
			biasanya.
		\end{eulercomment}
		\begin{eulerprompt}
			>shrinkwindow(>smaller);  ...
			>mosaicplot(BT',srows=YT,scols=P,color=CP,style="#"); ...
			>shrinkwindow():
		\end{eulerprompt}
		\eulerimg{25}{images/Tugas 7-024.png}
		\begin{eulercomment}
			Kita juga dapat membuat diagram lingkaran (pie chart). Karena hitam dan kuning membentuk koalisi,
			kita mengurutkan ulang elemen-elemennya.
		\end{eulercomment}
		\begin{eulerprompt}
			>i=[1,3,5,4,2]; piechart(BW[6,3:7][i],color=CP[i],lab=P[i]):
		\end{eulerprompt}
		\eulerimg{25}{images/Tugas 7-025.png}
		\begin{eulercomment}
			Berikut adalah jenis plot lainnya.
		\end{eulercomment}
		\begin{eulerprompt}
			>starplot(normal(1,10)+4,lab=1:10,>rays):
		\end{eulerprompt}
		\eulerimg{25}{images/Tugas 7-026.png}
		\begin{eulercomment}
			Beberapa plot dalam \texttt{plot2d} sangat baik untuk statistika. Berikut adalah plot impuls dari data acak
			yang berdistribusi seragam di $[0,1]$.
		\end{eulercomment}
		\begin{eulerprompt}
			>plot2d(makeimpulse(1:10,random(1,10)),>bar):
		\end{eulerprompt}
		\eulerimg{25}{images/Tugas 7-027.png}
		\begin{eulercomment}
			Namun untuk data yang berdistribusi eksponensial, kita mungkin memerlukan plot logaritmik.
		\end{eulercomment}
		\begin{eulerprompt}
			>logimpulseplot(1:10,-log(random(1,10))*10):
		\end{eulerprompt}
		\eulerimg{25}{images/Tugas 7-028.png}
		\begin{eulercomment}
			Fungsi \texttt{columnsplot()} lebih mudah digunakan, karena hanya memerlukan vektor nilai. Selain itu, fungsi
			ini dapat mengatur labelnya ke apa pun yang kita inginkan, seperti yang sudah kita demonstrasikan
			dalam tutorial ini.
			
			Berikut adalah aplikasi lain, di mana kita menghitung karakter dalam sebuah kalimat dan memplot
			statistiknya.
		\end{eulercomment}
		\begin{eulerprompt}
			>v=strtochar("the quick brown fox jumps over the lazy dog"); ...
			>w=ascii("a"):ascii("z"); x=getmultiplicities(w,v); ...
			>cw=[]; for k=w; cw=cw|char(k); end; ...
			>columnsplot(x,lab=cw,width=0.05):
		\end{eulerprompt}
		\eulerimg{25}{images/Tugas 7-029.png}
		\begin{eulercomment}
			Dimungkinkan juga untuk mengatur sumbu secara manual.
		\end{eulercomment}
		\begin{eulerprompt}
			>n=10; p=0.4; i=0:n; x=bin(n,i)*p^i*(1-p)^(n-i); ...
			>columnsplot(x,lab=i,width=0.05,<frame,<grid); ...
			>yaxis(0,0:0.1:1,style="->",>left); xaxis(0,style="."); ...
			>label("p",0,0.25), label("i",11,0); ...
			>textbox(["Binomial distribution","with p=0.4"]):
		\end{eulerprompt}
		\eulerimg{25}{images/Tugas 7-030.png}
		\begin{eulercomment}
			Berikut adalah cara untuk memplot frekuensi angka dalam suatu vektor.
			
			Kita membuat vektor bilangan acak bulat 1 hingga 10.
		\end{eulercomment}
		\begin{eulerprompt}
			>v:=intrandom(1,10,10)
		\end{eulerprompt}
		\begin{euleroutput}
			[8,  5,  8,  8,  6,  8,  8,  3,  5,  5]
		\end{euleroutput}
		\begin{eulercomment}
			Kemudian mengekstrak angka-angka unik dalam \texttt{v}.
		\end{eulercomment}
		\begin{eulerprompt}
			>vu:=unique(v)
		\end{eulerprompt}
		\begin{euleroutput}
			[3,  5,  6,  8]
		\end{euleroutput}
		\begin{eulercomment}
			Dan memplot frekuensinya dalam plot kolom.
		\end{eulercomment}
		\begin{eulerprompt}
			>columnsplot(getmultiplicities(vu,v),lab=vu,style="/"):
		\end{eulerprompt}
		\eulerimg{25}{images/Tugas 7-031.png}
		\begin{eulercomment}
			Kita ingin mendemonstrasikan fungsi untuk distribusi empiris dari nilai-nilai.
		\end{eulercomment}
		\begin{eulerprompt}
			>x=normal(1,20);
		\end{eulerprompt}
		\begin{eulercomment}
			Fungsi \texttt{empdist(x,vs)} memerlukan larik (array) nilai yang terurut. Jadi kita harus mengurutkan $x$
			sebelum dapat menggunakannya.
		\end{eulercomment}
		\begin{eulerprompt}
			>xs=sort(x);
		\end{eulerprompt}
		\begin{eulercomment}
			Kemudian kita memplot distribusi empiris dan beberapa batang kepadatan ke dalam satu plot. Alih-alih
			plot batang untuk distribusinya, kali ini kita menggunakan plot gigi gergaji (sawtooth plot).
		\end{eulercomment}
		\begin{eulerprompt}
			>figure(2,1); ...
			>figure(1); plot2d("empdist",-4,4;xs); ...
			>figure(2); plot2d(histo(x,v=-4:0.2:4,<bar));  ...
			>figure(0):
		\end{eulerprompt}
		\eulerimg{25}{images/Tugas 7-032.png}
		\begin{eulercomment}
			Plot pencar (scatter plot) mudah dibuat di Euler dengan plot titik biasa. Grafik berikut menunjukkan
			bahwa $X$ dan $X+Y$ jelas berkorelasi positif.
		\end{eulercomment}
		\begin{eulerprompt}
			>x=normal(1,100); plot2d(x,x+rotright(x),>points,style=".."):
		\end{eulerprompt}
		\eulerimg{25}{images/Tugas 7-033.png}
		\begin{eulercomment}
			Sering kali, kita ingin membandingkan dua sampel dari distribusi yang berbeda. Hal ini dapat
			dilakukan dengan plot kuantil-kuantil (quantile-quantile-plot).
			
			Sebagai pengujian, kita mencoba distribusi student-$t$ dan distribusi normal.
		\end{eulercomment}
		\begin{eulerprompt}
			>x=randt(1,1000,5); y=randnormal(1,1000,mean(x),dev(x)); ...
			>plot2d("x",r=6,style="--",yl="normal",xl="student-t",>vertical); ...
			>plot2d(sort(x),sort(y),>points,color=red,style="x",>add):
		\end{eulerprompt}
		\eulerimg{25}{images/Tugas 7-034.png}
		\begin{eulercomment}
			Plot tersebut dengan jelas menunjukkan bahwa nilai-nilai yang berdistribusi normal cenderung lebih
			kecil di ujung-ujung ekstremnya.
			
			Jika kita memiliki dua distribusi dengan ukuran berbeda, kita dapat memperluas yang lebih kecil atau
			menyusutkan yang lebih besar. Fungsi berikut baik untuk keduanya. Fungsi ini mengambil nilai median
			dengan persentase antara 0 dan 1.
		\end{eulercomment}
		\begin{eulerprompt}
			>function medianexpand (x,n) := median(x,p=linspace(0,1,n-1));
		\end{eulerprompt}
		\begin{eulercomment}
			Mari kita bandingkan dua distribusi yang sama.
		\end{eulercomment}
		\begin{eulerprompt}
			>x=random(1000); y=random(400); ...
			>plot2d("x",0,1,style="--"); ...
			>plot2d(sort(medianexpand(x,400)),sort(y),>points,color=red,style="x",>add):
		\end{eulerprompt}
		\eulerimg{25}{images/Tugas 7-035.png}
		
		\eulerheading{Regresi dan Korelasi}
		\begin{eulercomment}
			Regresi linear dapat dilakukan dengan fungsi \texttt{polyfit()} atau berbagai fungsi pencocokan (fit) lainnya.
			
			Sebagai permulaan, kita mencari garis regresi untuk data univariat dengan \texttt{polyfit(x,y,1)}.
		\end{eulercomment}
		\begin{eulerprompt}
			>x=1:10; y=[2,3,1,5,6,3,7,8,9,8]; writetable(x'|y',labc=["x","y"])
		\end{eulerprompt}
		\begin{euleroutput}
			         x         y
			         1         2
			         2         3
			         3         1
			         4         5
			         5         6
			         6         3
			         7         7
			         8         8
			         9         9
			        10         8
		\end{euleroutput}
		\begin{eulercomment}
			Kita ingin membandingkan pencocokan tanpa bobot dan berbobot. Pertama, koefisien dari pencocokan
			linear.
		\end{eulercomment}
		\begin{eulerprompt}
			>p=polyfit(x,y,1)
		\end{eulerprompt}
		\begin{euleroutput}
			[0.733333,  0.812121]
		\end{euleroutput}
		\begin{eulercomment}
			Sekarang koefisien dengan bobot yang menekankan nilai-nilai terakhir.
		\end{eulercomment}
		\begin{eulerprompt}
			>w &= "exp(-(x-10)^2/10)"; pw=polyfit(x,y,1,w=w(x))
		\end{eulerprompt}
		\begin{euleroutput}
			[4.71566,  0.38319]
		\end{euleroutput}
		\begin{eulercomment}
			Kita memasukkan semuanya ke dalam satu plot untuk titik-titik dan garis regresinya, serta untuk bobot
			yang digunakan.
		\end{eulercomment}
		\begin{eulerprompt}
			>figure(2,1);  ...
			>figure(1); statplot(x,y,"b",xl="Regression"); ...
			>  plot2d("evalpoly(x,p)",>add,color=blue,style="--"); ...
			>  plot2d("evalpoly(x,pw)",5,10,>add,color=red,style="--"); ...
			>figure(2); plot2d(w,1,10,>filled,style="/",fillcolor=red,xl=w); ...
			>figure(0):
		\end{eulerprompt}
		\eulerimg{25}{images/Tugas 7-036.png}
		\begin{eulercomment}
			Untuk contoh lain, kita membaca data survei mahasiswa, usia mereka, usia orang tua mereka, dan jumlah
			saudara kandung dari sebuah file.
			
			Tabel ini berisi \texttt{"m"} dan \texttt{"f"} di kolom kedua. Kita menggunakan variabel \texttt{tok2} untuk mengatur terjemahan
			yang tepat alih-alih membiarkan \texttt{readtable()} mengumpulkan terjemahannya sendiri.
		\end{eulercomment}
		\begin{eulerprompt}
			>{MS,hd}:=readtable("table1.dat",tok2:=["m","f"]);  ...
			>writetable(MS,labc=hd,tok2:=["m","f"]);
		\end{eulerprompt}
		\begin{euleroutput}
			    Person       Sex       Age    Mother    Father  Siblings
			         1         m        29        58        61         1
			         2         f        26        53        54         2
			         3         m        24        49        55         1
			         4         f        25        56        63         3
			         5         f        25        49        53         0
			         6         f        23        55        55         2
			         7         m        23        48        54         2
			         8         m        27        56        58         1
			         9         m        25        57        59         1
			        10         m        24        50        54         1
			        11         f        26        61        65         1
			        12         m        24        50        52         1
			        13         m        29        54        56         1
			        14         m        28        48        51         2
			        15         f        23        52        52         1
			        16         m        24        45        57         1
			        17         f        24        59        63         0
			        18         f        23        52        55         1
			        19         m        24        54        61         2
			        20         f        23        54        55         1
		\end{euleroutput}
		\begin{eulercomment}
			Bagaimana usia-usia tersebut saling bergantung satu sama lain? Gambaran pertama diperoleh dari plot
			pencar berpasangan (pairwise scatterplot).
		\end{eulercomment}
		\begin{eulerprompt}
			>scatterplots(tablecol(MS,3:5),hd[3:5]):
		\end{eulerprompt}
		\eulerimg{25}{images/Tugas 7-037.png}
		\begin{eulercomment}
			Jelas terlihat bahwa usia ayah dan ibu saling bergantung satu sama lain. Mari kita tentukan dan plot
			garis regresinya.
		\end{eulercomment}
		\begin{eulerprompt}
			>cs:=MS[,4:5]'; ps:=polyfit(cs[1],cs[2],1)
		\end{eulerprompt}
		\begin{euleroutput}
			[17.3789,  0.740964]
		\end{euleroutput}
		\begin{eulercomment}
			Model ini jelas keliru. Garis regresinya akan menjadi $s=17+0.74t$, di mana $t$ adalah usia ibu dan $s$
			adalah usia ayah. Perbedaan usia mungkin sedikit bergantung pada usia, tetapi tidak sebanyak itu.
			
			Sebaliknya, kita menduga fungsi seperti $s=a+t$. Maka $a$ adalah rata-rata dari $s-t$. Nilai ini merupakan
			rata-rata selisih usia antara ayah dan ibu.
		\end{eulercomment}
		\begin{eulerprompt}
			>da:=mean(cs[2]-cs[1])
		\end{eulerprompt}
		\begin{euleroutput}
			3.65
		\end{euleroutput}
		\begin{eulercomment}
			Mari kita plot ini ke dalam satu plot pencar.
		\end{eulercomment}
		\begin{eulerprompt}
			>plot2d(cs[1],cs[2],>points);  ...
			>plot2d("evalpoly(x,ps)",color=red,style=".",>add);  ...
			>plot2d("x+da",color=blue,>add):
		\end{eulerprompt}
		\eulerimg{25}{images/Tugas 7-038.png}
		\begin{eulercomment}
			Berikut adalah plot kotak dari kedua usia tersebut. Ini hanya menunjukkan bahwa usia keduanya
			berbeda.
		\end{eulercomment}
		\begin{eulerprompt}
			>boxplot(cs,["mothers","fathers"]):
		\end{eulerprompt}
		\eulerimg{25}{images/Tugas 7-039.png}
		\begin{eulercomment}
			Menarik untuk dicatat bahwa selisih mediannya tidak sebesar selisih rata-ratanya.
		\end{eulercomment}
		\begin{eulerprompt}
			>median(cs[2])-median(cs[1])
		\end{eulerprompt}
		\begin{euleroutput}
			1.5
		\end{euleroutput}
		\begin{eulercomment}
			Koefisien korelasi menunjukkan adanya korelasi positif.
		\end{eulercomment}
		\begin{eulerprompt}
			>correl(cs[1],cs[2])
		\end{eulerprompt}
		\begin{euleroutput}
			0.7588307236
		\end{euleroutput}
		\begin{eulercomment}
			Korelasi peringkat (rank correlation) adalah ukuran untuk urutan yang sama pada kedua vektor.
			Nilainya juga cukup positif.
		\end{eulercomment}
		\begin{eulerprompt}
			>rankcorrel(cs[1],cs[2])
		\end{eulerprompt}
		\begin{euleroutput}
			0.758925292358
		\end{euleroutput}
		
		\eulerheading{Membuat Fungsi Baru}
		\begin{eulercomment}
			Tentu saja, bahasa EMT dapat digunakan untuk memprogram fungsi-fungsi baru. Misalnya, kita
			mendefinisikan fungsi kemencengan (skewness).
		\end{eulercomment}
		\begin{eulerformula}
			\[
			\text{sk}(x) = \dfrac{\sqrt{n} \sum_i (x_i-m)^3}{\left(\sum_i (x_i-m)^2\right)^{3/2}}
			\]
		\end{eulerformula}
		\begin{eulercomment}
			di mana $m$ adalah rata-rata dari $x$.
		\end{eulercomment}
		\begin{eulerprompt}
			>function skew (x:vector) ...
		\end{eulerprompt}
		\begin{eulerudf}
			m=mean(x);
			return sqrt(cols(x))*sum((x-m)^3)/(sum((x-m)^2))^(3/2);
			endfunction
		\end{eulerudf}
		\begin{eulercomment}
			Seperti yang Anda lihat, kita dapat dengan mudah menggunakan bahasa matriks untuk mendapatkan
			implementasi yang sangat singkat dan efisien. Mari kita coba fungsi ini.
		\end{eulercomment}
		\begin{eulerprompt}
			>data=normal(20); skew(normal(10))
		\end{eulerprompt}
		\begin{euleroutput}
			-0.198710316203
		\end{euleroutput}
		\begin{eulercomment}
			Berikut adalah fungsi lain, yang disebut koefisien kemencengan Pearson.
		\end{eulercomment}
		\begin{eulerprompt}
			>function skew1 (x) := 3*(mean(x)-median(x))/dev(x)
			>skew1(data)
		\end{eulerprompt}
		\begin{euleroutput}
			-0.0801873249135
		\end{euleroutput}
		
		\eulerheading{Simulasi Monte Carlo}
		\begin{eulercomment}
			Euler dapat digunakan untuk mensimulasikan kejadian acak. Kita telah melihat contoh sederhana di
			atas. Berikut adalah contoh lainnya, yang mensimulasikan 1000 kali lemparan 3 dadu, dan mencari
			distribusi dari jumlah mata dadunya.
		\end{eulercomment}
		\begin{eulerprompt}
			>ds:=sum(intrandom(1000,3,6))';  fs=getmultiplicities(3:18,ds)
		\end{eulerprompt}
		\begin{euleroutput}
			[5,  17,  35,  44,  75,  97,  114,  116,  143,  116,  104,  53,  40,  22,  13,  6]
		\end{euleroutput}
		\begin{eulercomment}
			Kita dapat memplotnya sekarang.
		\end{eulercomment}
		\begin{eulerprompt}
			>columnsplot(fs,lab=3:18):
		\end{eulerprompt}
		\eulerimg{25}{images/Tugas 7-041.png}
		\begin{eulercomment}
			Menentukan distribusi yang diharapkan tidaklah begitu mudah. Kita menggunakan rekursi tingkat lanjut
			untuk ini.
			
			Fungsi berikut menghitung banyaknya cara angka $k$ dapat direpresentasikan sebagai penjumlahan dari $n$
			bilangan dalam rentang 1 hingga $m$. Fungsi ini bekerja secara rekursif dengan cara yang jelas.
		\end{eulercomment}
		\begin{eulerprompt}
			>function map countways (k; n, m) ...
		\end{eulerprompt}
		\begin{eulerudf}
			if n==1 then return k>=1 && k<=m
			else
			  sum=0; 
			  loop 1 to m; sum=sum+countways(k-#,n-1,m); end;
			  return sum;
			end;
			endfunction
		\end{eulerudf}
		\begin{eulercomment}
			Berikut adalah hasil untuk lima lemparan dadu bersisi lima dan tiga lemparan dadu bersisi enam.
		\end{eulercomment}
		\begin{eulerprompt}
			>countways(5:25,5,5)
		\end{eulerprompt}
		\begin{euleroutput}
			[1,  5,  15,  35,  70,  121,  185,  255,  320,  365,  381,  365,  320,  255,  185,  121,  70,  35,
			15,  5,  1]
		\end{euleroutput}
		\begin{eulerprompt}
			>cw=countways(3:18,3,6)
		\end{eulerprompt}
		\begin{euleroutput}
			[1,  3,  6,  10,  15,  21,  25,  27,  27,  25,  21,  15,  10,  6,  3,  1]
		\end{euleroutput}
		\begin{eulercomment}
			Kita menambahkan nilai harapan ke dalam plot.
		\end{eulercomment}
		\begin{eulerprompt}
			>plot2d(cw/6^3*1000,>add); plot2d(cw/6^3*1000,>points,>add):
		\end{eulerprompt}
		\eulerimg{25}{images/Tugas 7-042.png}
		\begin{eulercomment}
			Untuk simulasi lain, simpangan nilai rata-rata dari $n$ variabel acak berdistribusi normal $0\text{--}1$ adalah
			$1/\sqrt{n}$.
		\end{eulercomment}
		\begin{eulerprompt}
			>longformat; 1/sqrt(10)
		\end{eulerprompt}
		\begin{euleroutput}
			0.316227766017
		\end{euleroutput}
		\begin{eulercomment}
			Mari kita periksa hal ini dengan simulasi. Kita menghasilkan 10000 kali 10 vektor acak.
		\end{eulercomment}
		\begin{eulerprompt}
			>M=normal(10000,10); dev(mean(M)')
		\end{eulerprompt}
		\begin{euleroutput}
			0.319493614817
		\end{euleroutput}
		\begin{eulerprompt}
			>plot2d(mean(M)',>distribution):
		\end{eulerprompt}
		\eulerimg{25}{images/Tugas 7-043.png}
		\begin{eulercomment}
			Median dari 10 bilangan acak berdistribusi normal $0\text{--}1$ memiliki simpangan yang lebih besar.
		\end{eulercomment}
		\begin{eulerprompt}
			>dev(median(M)')
		\end{eulerprompt}
		\begin{euleroutput}
			0.374460271535
		\end{euleroutput}
		\begin{eulercomment}
			Karena kita dapat dengan mudah menghasilkan jalan acak (random walk), kita dapat mensimulasikan
			proses Wiener. Kita mengambil 1000 langkah dari 1000 proses. Kemudian kita memplot simpangan baku dan
			rata-rata dari langkah ke-$n$ proses-proses ini bersama dengan nilai harapannya dalam warna merah.
		\end{eulercomment}
		\begin{eulerprompt}
			>n=1000; m=1000; M=cumsum(normal(n,m)/sqrt(m)); ...
			>t=(1:n)/n; figure(2,1); ...
			>figure(1); plot2d(t,mean(M')'); plot2d(t,0,color=red,>add); ...
			>figure(2); plot2d(t,dev(M')'); plot2d(t,sqrt(t),color=red,>add); ...
			>figure(0):
		\end{eulerprompt}
		\eulerimg{25}{images/Tugas 7-044.png}
		
		\eulerheading{Pengujian (Uji Statistik)}
		\begin{eulercomment}
			Pengujian adalah alat penting dalam statistika. Di Euler, banyak uji yang telah diimplementasikan.
			Semua uji ini mengembalikan galat ($p$-value) yang kita terima jika kita menolak hipotesis nol.
			
			Sebagai contoh, kita menguji lemparan dadu untuk distribusi seragam. Pada 600 lemparan, kita
			mendapatkan nilai-nilai berikut, yang kita masukkan ke dalam uji chi-kuadrat.
		\end{eulercomment}
		\begin{eulerprompt}
			>chitest([90,103,114,101,103,89],dup(100,6)')
		\end{eulerprompt}
		\begin{euleroutput}
			0.498830517952
		\end{euleroutput}
		\begin{eulercomment}
			Uji chi-kuadrat juga memiliki mode yang menggunakan simulasi Monte Carlo untuk menguji statistik.
			Hasilnya seharusnya hampir sama. Parameter \texttt{>p} menafsirkan vektor $y$ sebagai vektor probabilitas.
		\end{eulercomment}
		\begin{eulerprompt}
			>chitest([90,103,114,101,103,89],dup(1/6,6)',>p,>montecarlo)
		\end{eulerprompt}
		\begin{euleroutput}
			0.526
		\end{euleroutput}
		\begin{eulercomment}
			Galat ini jauh terlalu besar. Jadi kita tidak dapat menolak distribusi seragam. Hal ini tidak
			membuktikan bahwa dadu kita adil, tetapi kita tidak dapat menolak hipotesis kita.
			
			Selanjutnya kita membangkitkan 1000 lemparan dadu menggunakan pembangkit bilangan acak, dan melakukan
			pengujian yang sama.
		\end{eulercomment}
		\begin{eulerprompt}
			>n=1000; t=random([1,n*6]); chitest(count(t*6,6),dup(n,6)')
		\end{eulerprompt}
		\begin{euleroutput}
			0.528028118442
		\end{euleroutput}
		\begin{eulercomment}
			Mari kita uji nilai rata-rata 200 dengan uji-$t$ ($t$-test).
		\end{eulercomment}
		\begin{eulerprompt}
			>s=200+normal([1,100])*10; ...
			>ttest(mean(s),dev(s),100,200)
		\end{eulerprompt}
		\begin{euleroutput}
			0.0218365848476
		\end{euleroutput}
		\begin{eulercomment}
			Fungsi \texttt{ttest()} memerlukan nilai rata-rata, simpangan, jumlah data, dan nilai rata-rata yang akan
			diuji.
			
			Sekarang mari kita periksa dua pengukuran untuk rata-rata yang sama. Kita menolak hipotesis bahwa
			keduanya memiliki rata-rata yang sama jika hasilnya $< 0.05$.
		\end{eulercomment}
		\begin{eulerprompt}
			>tcomparedata(normal(1,10),normal(1,10))
		\end{eulerprompt}
		\begin{euleroutput}
			0.38722000942
		\end{euleroutput}
		\begin{eulercomment}
			Jika kita menambahkan bias pada salah satu distribusi, kita mendapatkan lebih banyak penolakan.
			Ulangi simulasi ini beberapa kali untuk melihat efeknya.
		\end{eulercomment}
		\begin{eulerprompt}
			>tcomparedata(normal(1,10),normal(1,10)+2)
		\end{eulerprompt}
		\begin{euleroutput}
			5.60009101758e-07
		\end{euleroutput}
		\begin{eulercomment}
			Pada contoh berikutnya, kita membangkitkan 20 lemparan dadu acak sebanyak 100 kali dan menghitung
			kemunculan angka satu di dalamnya. Rata-ratanya haruslah $20/6 \approx 3.3$ angka satu.
		\end{eulercomment}
		\begin{eulerprompt}
			>R=random(100,20); R=sum(R*6<=1)'; mean(R)
		\end{eulerprompt}
		\begin{euleroutput}
			3.28
		\end{euleroutput}
		\begin{eulercomment}
			Sekarang kita membandingkan jumlah angka satu tersebut dengan distribusi binomial. Pertama kita
			memplot distribusi angka satu.
		\end{eulercomment}
		\begin{eulerprompt}
			>plot2d(R,distribution=max(R)+1,even=1,style="\/"):
		\end{eulerprompt}
		\eulerimg{25}{images/Tugas 7-045.png}
		\begin{eulerprompt}
			>t=count(R,21);
		\end{eulerprompt}
		\begin{eulercomment}
			Kemudian kita menghitung nilai harapannya.
		\end{eulercomment}
		\begin{eulerprompt}
			>n=0:20; b=bin(20,n)*(1/6)^n*(5/6)^(20-n)*100;
		\end{eulerprompt}
		\begin{eulercomment}
			Kita harus mengumpulkan beberapa angka untuk mendapatkan kategori yang cukup besar.
		\end{eulercomment}
		\begin{eulerprompt}
			>t1=sum(t[1:2])|t[3:7]|sum(t[8:21]); ...
			>b1=sum(b[1:2])|b[3:7]|sum(b[8:21]);
		\end{eulerprompt}
		\begin{eulercomment}
			Uji chi-kuadrat menolak hipotesis bahwa distribusi kita adalah distribusi binomial jika hasilnya
			$< 0.05$.
		\end{eulercomment}
		\begin{eulerprompt}
			>chitest(t1,b1)
		\end{eulerprompt}
		\begin{euleroutput}
			0.53921579764
		\end{euleroutput}
		\begin{eulercomment}
			Contoh berikut berisi hasil dari dua kelompok orang (katakanlah pria dan wanita) yang memilih salah
			satu dari enam partai.
		\end{eulercomment}
		\begin{eulerprompt}
			>A=[23,37,43,52,64,74;27,39,41,49,63,76];  ...
			>  writetable(A,wc=6,labr=["m","f"],labc=1:6)
		\end{eulerprompt}
		\begin{euleroutput}
			           1     2     3     4     5     6
			     m    23    37    43    52    64    74
			     f    27    39    41    49    63    76
		\end{euleroutput}
		\begin{eulercomment}
			Kita ingin menguji independensi pilihan suara terhadap jenis kelamin. Uji tabel $\chi^2$ melakukan hal
			ini. Hasilnya terlalu besar untuk menolak independensi. Jadi kita tidak dapat mengatakan apakah
			pilihan suara bergantung pada jenis kelamin berdasarkan data ini.
		\end{eulercomment}
		\begin{eulerprompt}
			>tabletest(A)
		\end{eulerprompt}
		\begin{euleroutput}
			0.990701632326
		\end{euleroutput}
		\begin{eulercomment}
			Berikut adalah tabel yang diharapkan, jika kita mengasumsikan frekuensi pemilihan yang diamati.
		\end{eulercomment}
		\begin{eulerprompt}
			>writetable(expectedtable(A),wc=6,dc=1,labr=["m","f"],labc=1:6)
		\end{eulerprompt}
		\begin{euleroutput}
			           1     2     3     4     5     6
			     m  24.9  37.9  41.9  50.3  63.3  74.7
			     f  25.1  38.1  42.1  50.7  63.7  75.3
		\end{euleroutput}
		\begin{eulercomment}
			Kita dapat menghitung koefisien kontingensi terkoreksi. Karena nilainya sangat dekat dengan 0, kita
			menyimpulkan bahwa pilihan suara tidak bergantung pada jenis kelamin.
		\end{eulercomment}
		\begin{eulerprompt}
			>contingency(A)
		\end{eulerprompt}
		\begin{euleroutput}
			0.0427225484717
		\end{euleroutput}
		
		\eulerheading{Beberapa Uji Lainnya}
		\begin{eulercomment}
			Selanjutnya kita menggunakan analisis varians (uji-$F$) untuk menguji tiga sampel data berdistribusi
			normal terhadap nilai rata-rata yang sama. Metode ini disebut ANOVA (analysis of variance). Di Euler,
			fungsi \texttt{varanalysis()} digunakan.
		\end{eulercomment}
		\begin{eulerprompt}
			>x1=[109,111,98,119,91,118,109,99,115,109,94]; mean(x1),
		\end{eulerprompt}
		\begin{euleroutput}
			106.545454545
		\end{euleroutput}
		\begin{eulerprompt}
			>x2=[120,124,115,139,114,110,113,120,117]; mean(x2),
		\end{eulerprompt}
		\begin{euleroutput}
			119.111111111
		\end{euleroutput}
		\begin{eulerprompt}
			>x3=[120,112,115,110,105,134,105,130,121,111]; mean(x3)
		\end{eulerprompt}
		\begin{euleroutput}
			116.3
		\end{euleroutput}
		\begin{eulerprompt}
			>varanalysis(x1,x2,x3)
		\end{eulerprompt}
		\begin{euleroutput}
			0.0138048221371
		\end{euleroutput}
		\begin{eulercomment}
			Ini berarti kita menolak hipotesis nilai rata-rata yang sama. Kita melakukannya dengan probabilitas
			galat sebesar 1,3\%.
			
			Terdapat juga uji median, yang menolak sampel data dengan distribusi rata-rata yang berbeda dengan
			menguji median dari sampel gabungan.
		\end{eulercomment}
		\begin{eulerprompt}
			>a=[56,66,68,49,61,53,45,58,54];
			>b=[72,81,51,73,69,78,59,67,65,71,68,71];
			>mediantest(a,b)
		\end{eulerprompt}
		\begin{euleroutput}
			0.0241724220052
		\end{euleroutput}
		\begin{eulercomment}
			Uji kesamaan lainnya adalah uji peringkat (rank test). Uji ini jauh lebih tajam daripada uji median.
		\end{eulercomment}
		\begin{eulerprompt}
			>ranktest(a,b)
		\end{eulerprompt}
		\begin{euleroutput}
			0.00199969612469
		\end{euleroutput}
		\begin{eulercomment}
			Pada contoh berikut, kedua distribusi memiliki rata-rata yang sama.
		\end{eulercomment}
		\begin{eulerprompt}
			>ranktest(random(1,100),random(1,50)*3-1)
		\end{eulerprompt}
		\begin{euleroutput}
			0.129608141484
		\end{euleroutput}
		\begin{eulercomment}
			Mari kita coba mensimulasikan dua perlakuan $a$ dan $b$ yang diterapkan pada orang yang berbeda.
		\end{eulercomment}
		\begin{eulerprompt}
			>a=[8.0,7.4,5.9,9.4,8.6,8.2,7.6,8.1,6.2,8.9];
			>b=[6.8,7.1,6.8,8.3,7.9,7.2,7.4,6.8,6.8,8.1];
		\end{eulerprompt}
		\begin{eulercomment}
			Uji tanda (signum test) memutuskan apakah $a$ lebih baik daripada $b$.
		\end{eulercomment}
		\begin{eulerprompt}
			>signtest(a,b)
		\end{eulerprompt}
		\begin{euleroutput}
			0.0546875
		\end{euleroutput}
		\begin{eulercomment}
			Galat ini terlalu besar. Kita tidak dapat menolak bahwa $a$ sama baiknya dengan $b$.
			
			Uji Wilcoxon lebih tajam daripada uji ini, tetapi bergantung pada nilai kuantitatif dari selisihnya.
		\end{eulercomment}
		\begin{eulerprompt}
			>wilcoxon(a,b)
		\end{eulerprompt}
		\begin{euleroutput}
			0.0296680599405
		\end{euleroutput}
		\begin{eulercomment}
			Mari kita coba dua uji lagi menggunakan deret yang dibangkitkan.
		\end{eulercomment}
		\begin{eulerprompt}
			>wilcoxon(normal(1,20),normal(1,20)-1)
		\end{eulerprompt}
		\begin{euleroutput}
			0.0068706451766
		\end{euleroutput}
		\begin{eulerprompt}
			>wilcoxon(normal(1,20),normal(1,20))
		\end{eulerprompt}
		\begin{euleroutput}
			0.275145971064
		\end{euleroutput}
		
		\eulerheading{Bilangan Acak}
		\begin{eulercomment}
			Berikut adalah pengujian untuk pembangkit bilangan acak. Euler menggunakan pembangkit yang sangat
			baik, sehingga kita tidak perlu mengkhawatirkan adanya masalah.
			
			Pertama kita membangkitkan sepuluh juta bilangan acak dalam $[0,1]$.
		\end{eulercomment}
		\begin{eulerprompt}
			>n:=10000000; r:=random(1,n);
		\end{eulerprompt}
		\begin{eulercomment}
			Selanjutnya kita menghitung jarak antara dua bilangan yang kurang dari 0.05.
		\end{eulercomment}
		\begin{eulerprompt}
			>a:=0.05; d:=differences(nonzeros(r<a));
		\end{eulerprompt}
		\begin{eulercomment}
			Terakhir, kita memplot berapa kali setiap jarak muncul, dan membandingkannya dengan nilai harapan.
		\end{eulercomment}
		\begin{eulerprompt}
			>m=getmultiplicities(1:100,d); plot2d(m); ...
			>  plot2d("n*(1-a)^(x-1)*a^2",color=red,>add):
		\end{eulerprompt}
		\eulerimg{25}{images/Tugas 7-046.png}
		\begin{eulercomment}
			Hapus data.
		\end{eulercomment}
		\begin{eulerprompt}
			>remvalue n;
		\end{eulerprompt}
		
		\eulerheading{Pengantar untuk Pengguna Proyek R}
		\begin{eulercomment}
			Jelas bahwa EMT tidak bersaing dengan R sebagai paket statistika. Namun, terdapat banyak prosedur dan
			fungsi statistika yang tersedia di EMT juga. Jadi EMT dapat memenuhi kebutuhan dasar. Bagaimanapun
			juga, EMT dilengkapi dengan paket numerik dan sistem aljabar komputer.
			
			Notebook ini ditujukan untuk Anda jika Anda sudah terbiasa dengan R, tetapi perlu mengetahui
			perbedaan sintaks antara EMT dan R. Kami mencoba memberikan gambaran umum tentang hal-hal yang jelas
			maupun yang kurang jelas yang perlu Anda ketahui.
			
			Selain itu, kita melihat cara-cara untuk bertukar data antara kedua sistem tersebut.
			
			Perhatikan bahwa bagian ini masih dalam pengembangan.
		\end{eulercomment}
		
		\eulerheading{Sintaks Dasar}
		\begin{eulercomment}
			Hal pertama yang Anda pelajari di R adalah membuat vektor. Di EMT, perbedaan utamanya adalah operator
			\texttt{:} dapat menerima ukuran langkah (step size). Selain itu, operator ini memiliki daya ikat yang rendah.
		\end{eulercomment}
		\begin{eulerprompt}
			>n=10; 0:n/20:n-1
		\end{eulerprompt}
		\begin{euleroutput}
			[0,  0.5,  1,  1.5,  2,  2.5,  3,  3.5,  4,  4.5,  5,  5.5,  6,  6.5,  7,  7.5,  8,  8.5,  9]
		\end{euleroutput}
		\begin{eulercomment}
			Fungsi \texttt{c()} tidak ada. Dimungkinkan untuk menggunakan vektor guna menggabungkan elemen-elemen.
			
			Contoh berikut, seperti banyak contoh lainnya, berasal dari ``Introduction to R'' yang disertakan
			dengan proyek R. Jika Anda membaca PDF tersebut, Anda akan menemukan bahwa saya mengikuti alurnya
			dalam tutorial ini.
		\end{eulercomment}
		\begin{eulerprompt}
			>x=[10.4, 5.6, 3.1, 6.4, 21.7]; [x,0,x]
		\end{eulerprompt}
		\begin{euleroutput}
			[10.4,  5.6,  3.1,  6.4,  21.7,  0,  10.4,  5.6,  3.1,  6.4,  21.7]
		\end{euleroutput}
		\begin{eulercomment}
			Operator titik dua dengan ukuran langkah pada EMT digantikan oleh fungsi \texttt{seq()} di R. Kita dapat
			menulis fungsi ini di EMT.
		\end{eulercomment}
		\begin{eulerprompt}
			>function seq(a,b,c) := a:b:c; ...
			>seq(0,-0.1,-1)
		\end{eulerprompt}
		\begin{euleroutput}
			[0,  -0.1,  -0.2,  -0.3,  -0.4,  -0.5,  -0.6,  -0.7,  -0.8,  -0.9,  -1]
		\end{euleroutput}
		\begin{eulercomment}
			Fungsi \texttt{rep()} dari R tidak ada di EMT. Untuk masukan vektor, fungsi tersebut dapat ditulis sebagai
			berikut.
		\end{eulercomment}
		\begin{eulerprompt}
			>function rep(x:vector,n:index) := flatten(dup(x,n)); ...
			>rep(x,2)
		\end{eulerprompt}
		\begin{euleroutput}
			[10.4,  5.6,  3.1,  6.4,  21.7,  10.4,  5.6,  3.1,  6.4,  21.7]
		\end{euleroutput}
		\begin{eulercomment}
			Perhatikan bahwa \texttt{=} atau \texttt{:=} digunakan untuk penugasan (assignment). Operator \texttt{->} digunakan untuk
			satuan di EMT.
		\end{eulercomment}
		\begin{eulerprompt}
			>125km -> " miles"
		\end{eulerprompt}
		\begin{euleroutput}
			77.6713990297 miles
		\end{euleroutput}
		\begin{eulercomment}
			Operator \texttt{<-} untuk penugasan bagaimanapun juga menyesatkan, dan bukan ide yang baik dari R. Perintah
			berikut akan membandingkan $a$ dan $-4$ di EMT.
		\end{eulercomment}
		\begin{eulerprompt}
			>a=2; a<-4
		\end{eulerprompt}
		\begin{euleroutput}
			0
		\end{euleroutput}
		\begin{eulercomment}
			Di R, \texttt{a<-4<3} berfungsi, tetapi \texttt{a<-4<-3} tidak. Saya juga mengalami ambiguitas serupa di EMT,
			tetapi mencoba menghilangkannya sedikit demi sedikit.
			
			EMT dan R memiliki vektor bertipe boolean. Namun di EMT, angka 0 dan 1 digunakan untuk
			merepresentasikan false dan true. Di R, nilai true dan false tetap dapat digunakan dalam aritmetika
			biasa sama seperti di EMT.
		\end{eulercomment}
		\begin{eulerprompt}
			>x<5, %*x
		\end{eulerprompt}
		\begin{euleroutput}
			[0,  0,  1,  0,  0]
			[0,  0,  3.1,  0,  0]
		\end{euleroutput}
		\begin{eulercomment}
			EMT memunculkan pesan kesalahan (error) atau menghasilkan \texttt{NAN} bergantung pada flag \texttt{errors}.
		\end{eulercomment}
		\begin{eulerprompt}
			>errors off; 0/0, isNAN(sqrt(-1)), errors on;
		\end{eulerprompt}
		\begin{euleroutput}
			NAN
			1
		\end{euleroutput}
		\begin{eulercomment}
			String sama saja di R dan EMT. Keduanya menggunakan lokal (locale) saat ini, bukan dalam Unicode.
			
			Di R terdapat paket untuk Unicode. Di EMT, string dapat berupa string Unicode. String Unicode dapat
			diterjemahkan ke pengkodean lokal dan sebaliknya. Selain itu, \texttt{u"..."} dapat berisi entitas HTML.
		\end{eulercomment}
		\begin{eulerprompt}
			>u"&#169; Ren&eacut; Grothmann"
		\end{eulerprompt}
		\begin{euleroutput}
			© René Grothmann
		\end{euleroutput}
		\begin{eulercomment}
			Hal berikut mungkin tampil atau tidak tampil dengan benar di sistem Anda sebagai huruf A dengan titik
			dan garis di atasnya. Hal ini bergantung pada font yang Anda gunakan.
		\end{eulercomment}
		\begin{eulerprompt}
			>chartoutf([480])
		\end{eulerprompt}
		\begin{euleroutput}
			Ǡ
		\end{euleroutput}
		\begin{eulercomment}
			Penggabungan string dilakukan dengan \texttt{+} atau \texttt{|}. Ini dapat menyertakan angka, yang akan dicetak
			dalam format saat ini.
		\end{eulercomment}
		\begin{eulerprompt}
			>"pi = "+pi
		\end{eulerprompt}
		\begin{euleroutput}
			pi = 3.14159265359
		\end{euleroutput}
		
		\eulerheading{Pengindeksan}
		\begin{eulercomment}
			Sebagian besar waktu, hal ini akan bekerja seperti di R.
			
			Namun EMT akan menafsirkan indeks negatif dari bagian belakang vektor, sedangkan R menafsirkan \texttt{x[-n]}
			sebagai \texttt{x} tanpa elemen ke-$n$.
		\end{eulercomment}
		\begin{eulerprompt}
			>x, x[1:3], x[-2]
		\end{eulerprompt}
		\begin{euleroutput}
			[10.4,  5.6,  3.1,  6.4,  21.7]
			[10.4,  5.6,  3.1]
			6.4
		\end{euleroutput}
		\begin{eulercomment}
			Perilaku R dapat dicapai di EMT dengan \texttt{drop()}.
		\end{eulercomment}
		\begin{eulerprompt}
			>drop(x,2)
		\end{eulerprompt}
		\begin{euleroutput}
			[10.4,  3.1,  6.4,  21.7]
		\end{euleroutput}
		\begin{eulercomment}
			Vektor logika tidak diperlakukan secara berbeda sebagai indeks di EMT, berbeda dengan R. Anda perlu
			mengekstrak elemen bukan nol terlebih dahulu di EMT.
		\end{eulercomment}
		\begin{eulerprompt}
			>x, x>5, x[nonzeros(x>5)]
		\end{eulerprompt}
		\begin{euleroutput}
			[10.4,  5.6,  3.1,  6.4,  21.7]
			[1,  1,  0,  1,  1]
			[10.4,  5.6,  6.4,  21.7]
		\end{euleroutput}
		\begin{eulercomment}
			Sama seperti di R, vektor indeks dapat berisi pengulangan.
		\end{eulercomment}
		\begin{eulerprompt}
			>x[[1,2,2,1]]
		\end{eulerprompt}
		\begin{euleroutput}
			[10.4,  5.6,  5.6,  10.4]
		\end{euleroutput}
		\begin{eulercomment}
			Namun penamaan untuk indeks tidak dimungkinkan di EMT. Untuk sebuah paket statistika, hal ini sering
			kali diperlukan guna mempermudah akses ke elemen vektor.
			
			Untuk meniru perilaku ini, kita dapat mendefinisikan fungsi seperti berikut.
		\end{eulercomment}
		\begin{eulerprompt}
			>function mysel (v,i,s) := v[indexof(s,i)]; ...
			>s=["first","second","third","fourth"]; mysel(x,["first","third"],s)
		\end{eulerprompt}
		\begin{euleroutput}
			[10.4,  3.1]
		\end{euleroutput}
		
		\eulerheading{Tipe Data}
		\begin{eulercomment}
			EMT memiliki tipe data yang lebih tetap dibandingkan R. Jelas bahwa di R terdapat vektor yang dapat
			bertumbuh. Anda dapat mengatur vektor numerik kosong \texttt{v} dan menetapkan nilai pada elemen \texttt{v[17]}. Hal
			ini tidak dimungkinkan di EMT.
			
			Hal berikut ini agak kurang efisien.
		\end{eulercomment}
		\begin{eulerprompt}
			>v=[]; for i=1 to 10000; v=v|i; end;
		\end{eulerprompt}
		\begin{eulercomment}
			EMT sekarang akan menyusun vektor dengan \texttt{v} dan \texttt{i} yang ditambahkan pada tumpukan (stack) dan menyalin
			vektor tersebut kembali ke variabel global \texttt{v}.
			
			Cara yang lebih efisien adalah mendefinisikan vektor terlebih dahulu.
		\end{eulercomment}
		\begin{eulerprompt}
			>v=zeros(10000); for i=1 to 10000; v[i]=i; end;
		\end{eulerprompt}
		\begin{eulercomment}
			Untuk mengubah tipe data di EMT, Anda dapat menggunakan fungsi seperti \texttt{complex()}.
		\end{eulercomment}
		\begin{eulerprompt}
			>complex(1:4)
		\end{eulerprompt}
		\begin{euleroutput}
			[ 1+0i ,  2+0i ,  3+0i ,  4+0i  ]
		\end{euleroutput}
		\begin{eulercomment}
			Konversi ke string hanya dimungkinkan untuk tipe data dasar. Format saat ini digunakan untuk
			penggabungan string sederhana. Namun terdapat fungsi seperti \texttt{print()} atau \texttt{frac()}.
			
			Untuk vektor, Anda dapat dengan mudah menulis fungsi Anda sendiri.
		\end{eulercomment}
		\begin{eulerprompt}
			>function tostr (v) ...
		\end{eulerprompt}
		\begin{eulerudf}
			s="[";
			loop 1 to length(v);
			   s=s+print(v[#],2,0);
			   if #<length(v) then s=s+","; endif;
			end;
			return s+"]";
			endfunction
		\end{eulerudf}
		\begin{eulerprompt}
			>tostr(linspace(0,1,10))
		\end{eulerprompt}
		\begin{euleroutput}
			[0.00,0.10,0.20,0.30,0.40,0.50,0.60,0.70,0.80,0.90,1.00]
		\end{euleroutput}
		\begin{eulercomment}
			Untuk komunikasi dengan Maxima, terdapat fungsi \texttt{convertmxm()}, yang juga dapat digunakan untuk
			memformat vektor sebagai keluaran.
		\end{eulercomment}
		\begin{eulerprompt}
			>convertmxm(1:10)
		\end{eulerprompt}
		\begin{euleroutput}
			[1,2,3,4,5,6,7,8,9,10]
		\end{euleroutput}
		\begin{eulercomment}
			Untuk LaTeX, perintah \texttt{tex} dapat digunakan untuk memperoleh perintah LaTeX.
		\end{eulercomment}
		\begin{eulerprompt}
			>tex(&[1,2,3])
		\end{eulerprompt}
		\begin{euleroutput}
			\left[ 1 , 2 , 3 \right] 
		\end{euleroutput}
		
		\eulerheading{Faktor dan Tabel}
		\begin{eulercomment}
			Dalam pengantar R terdapat contoh dengan apa yang disebut faktor (factors).
			
			Berikut adalah daftar wilayah dari 30 negara bagian.
		\end{eulercomment}
		\begin{eulerprompt}
			>austates = ["tas", "sa", "qld", "nsw", "nsw", "nt", "wa", "wa", ...
			>"qld", "vic", "nsw", "vic", "qld", "qld", "sa", "tas", ...
			>"sa", "nt", "wa", "vic", "qld", "nsw", "nsw", "wa", ...
			>"sa", "act", "nsw", "vic", "vic", "act"];
		\end{eulerprompt}
		\begin{eulercomment}
			Misalkan kita memiliki pendapatan yang sesuai di setiap negara bagian.
		\end{eulercomment}
		\begin{eulerprompt}
			>incomes = [60, 49, 40, 61, 64, 60, 59, 54, 62, 69, 70, 42, 56, ...
			>61, 61, 61, 58, 51, 48, 65, 49, 49, 41, 48, 52, 46, ...
			>59, 46, 58, 43];
		\end{eulerprompt}
		\begin{eulercomment}
			Sekarang, kita ingin menghitung rata-rata pendapatan di wilayah-wilayah tersebut. Sebagai program
			statistika, R memiliki \texttt{factor()} dan \texttt{tapply()} untuk ini.
			
			EMT dapat melakukan ini dengan mencari indeks wilayah dalam daftar wilayah unik.
		\end{eulercomment}
		\begin{eulerprompt}
			>auterr=sort(unique(austates)); f=indexofsorted(auterr,austates)
		\end{eulerprompt}
		\begin{euleroutput}
			[6,  5,  4,  2,  2,  3,  8,  8,  4,  7,  2,  7,  4,  4,  5,  6,  5,  3,  8,  7,  4,  2,  2,  8,  5,
			1,  2,  7,  7,  1]
		\end{euleroutput}
		\begin{eulercomment}
			Pada titik tersebut, kita dapat menulis fungsi perulangan (loop) kita sendiri untuk melakukan hal-hal
			bagi satu faktor saja.
			
			Atau kita dapat meniru fungsi \texttt{tapply()} dengan cara berikut.
		\end{eulercomment}
		\begin{eulerprompt}
			>function map tappl (i; f$:call, cat, x) ...
		\end{eulerprompt}
		\begin{eulerudf}
			u=sort(unique(cat));
			f=indexof(u,cat);
			return f$(x[nonzeros(f==indexof(u,i))]);
			endfunction
		\end{eulerudf}
		\begin{eulercomment}
			Ini agak kurang efisien, karena menghitung wilayah unik untuk setiap $i$, tetapi cara ini berhasil.
		\end{eulercomment}
		\begin{eulerprompt}
			>tappl(auterr,"mean",austates,incomes)
		\end{eulerprompt}
		\begin{euleroutput}
			[44.5,  57.3333333333,  55.5,  53.6,  55,  60.5,  56,  52.25]
		\end{euleroutput}
		\begin{eulercomment}
			Perhatikan bahwa ini berfungsi untuk setiap vektor wilayah.
		\end{eulercomment}
		\begin{eulerprompt}
			>tappl(["act","nsw"],"mean",austates,incomes)
		\end{eulerprompt}
		\begin{euleroutput}
			[44.5,  57.3333333333]
		\end{euleroutput}
		\begin{eulercomment}
			Sekarang, paket statistika EMT mendefinisikan tabel sama seperti di R. Fungsi \texttt{readtable()} dan
			\texttt{writetable()} dapat digunakan untuk masukan dan keluaran.
			
			Jadi kita dapat mencetak rata-rata pendapatan negara bagian di wilayah-wilayah tersebut dengan cara
			yang mudah dibaca.
		\end{eulercomment}
		\begin{eulerprompt}
			>writetable(tappl(auterr,"mean",austates,incomes),labc=auterr,wc=7)
		\end{eulerprompt}
		\begin{euleroutput}
			    act    nsw     nt    qld     sa    tas    vic     wa
			   44.5  57.33   55.5   53.6     55   60.5     56  52.25
		\end{euleroutput}
		\begin{eulercomment}
			Kita juga dapat mencoba meniru perilaku R sepenuhnya.
			
			Faktor-faktor tersebut jelas harus disimpan dalam sebuah koleksi (collection) bersama tipe dan
			kategorinya (negara bagian dan wilayah dalam contoh kita). Untuk EMT, kita menambahkan indeks yang
			telah dihitung sebelumnya.
		\end{eulercomment}
		\begin{eulerprompt}
			>function makef (t) ...
		\end{eulerprompt}
		\begin{eulerudf}
			## Data faktor
			## Mengembalikan koleksi dengan data t, data unik, indeks.
			## Lihat: tapply
			u=sort(unique(t));
			return {{t,u,indexofsorted(u,t)}};
			endfunction
		\end{eulerudf}
		\begin{eulerprompt}
			>statef=makef(austates);
		\end{eulerprompt}
		\begin{eulercomment}
			Sekarang elemen ketiga dari koleksi akan berisi indeksnya.
		\end{eulercomment}
		\begin{eulerprompt}
			>statef[3]
		\end{eulerprompt}
		\begin{euleroutput}
			[6,  5,  4,  2,  2,  3,  8,  8,  4,  7,  2,  7,  4,  4,  5,  6,  5,  3,  8,  7,  4,  2,  2,  8,  5,
			1,  2,  7,  7,  1]
		\end{euleroutput}
		\begin{eulercomment}
			Sekarang kita dapat meniru \texttt{tapply()} dengan cara berikut. Ini akan mengembalikan tabel sebagai koleksi
			data tabel dan judul kolom.
		\end{eulercomment}
		\begin{eulerprompt}
			>function tapply (t:vector,tf,f$:call) ...
		\end{eulerprompt}
		\begin{eulerudf}
			## Membuat tabel dari data dan faktor
			## tf : keluaran dari makef()
			## Lihat: makef
			uf=tf[2]; f=tf[3]; x=zeros(length(uf));
			for i=1 to length(uf);
			   ind=nonzeros(f==i);
			   if length(ind)==0 then x[i]=NAN;
			   else x[i]=f$(t[ind]);
			   endif;
			end;
			return {{x,uf}};
			endfunction
		\end{eulerudf}
		\begin{eulercomment}
			Kita tidak menambahkan banyak pemeriksaan tipe di sini. Satu-satunya tindakan pencegahan menyangkut
			kategori (faktor) tanpa data. Namun sebaiknya tetap dilakukan pemeriksaan terhadap panjang $t$ yang
			benar dan kebenaran koleksi \texttt{tf}.
			
			Tabel ini dapat dicetak sebagai tabel dengan \texttt{writetable()}.
		\end{eulercomment}
		\begin{eulerprompt}
			>writetable(tapply(incomes,statef,"mean"),wc=7)
		\end{eulerprompt}
		\begin{euleroutput}
			    act    nsw     nt    qld     sa    tas    vic     wa
			   44.5  57.33   55.5   53.6     55   60.5     56  52.25
		\end{euleroutput}
		
		\eulerheading{Array (Larik)}
		\begin{eulercomment}
			EMT hanya memiliki dua dimensi untuk larik (array). Tipe datanya disebut matriks. Namun, akan mudah
			untuk menulis fungsi untuk dimensi yang lebih tinggi atau pustaka C untuk hal ini.
			
			R memiliki lebih dari dua dimensi. Di R, larik adalah vektor dengan bidang dimensi.
			
			Di EMT, vektor adalah matriks dengan satu baris. Vektor dapat diubah menjadi matriks dengan \texttt{redim()}.
		\end{eulercomment}
		\begin{eulerprompt}
			>shortformat; X=redim(1:20,4,5)
		\end{eulerprompt}
		\begin{euleroutput}
			        1         2         3         4         5 
			        6         7         8         9        10 
			       11        12        13        14        15 
			       16        17        18        19        20 
		\end{euleroutput}
		\begin{eulercomment}
			Ekstraksi baris dan kolom, atau sub-matriks, sangat mirip seperti di R.
		\end{eulercomment}
		\begin{eulerprompt}
			>X[,2:3]
		\end{eulerprompt}
		\begin{euleroutput}
			        2         3 
			        7         8 
			       12        13 
			       17        18 
		\end{euleroutput}
		\begin{eulercomment}
			Namun, di R dimungkinkan untuk menetapkan daftar indeks tertentu dari vektor ke suatu nilai. Hal yang
			sama di EMT hanya dimungkinkan dengan perulangan (loop).
		\end{eulercomment}
		\begin{eulerprompt}
			>function setmatrixvalue (M, i, j, v) ...
		\end{eulerprompt}
		\begin{eulerudf}
			loop 1 to max(length(i),length(j),length(v))
			   M[i{#},j{#}] = v{#};
			end;
			endfunction
		\end{eulerudf}
		\begin{eulercomment}
			Kita mendemonstrasikan hal ini untuk menunjukkan bahwa matriks diteruskan berdasarkan referensi
			(passed by reference) di EMT. Jika Anda tidak ingin mengubah matriks asli \texttt{M}, Anda perlu menyalinnya
			di dalam fungsi.
		\end{eulercomment}
		\begin{eulerprompt}
			>setmatrixvalue(X,1:3,3:-1:1,0); X,
		\end{eulerprompt}
		\begin{euleroutput}
			        1         2         0         4         5 
			        6         0         8         9        10 
			        0        12        13        14        15 
			       16        17        18        19        20 
		\end{euleroutput}
		\begin{eulercomment}
			Hasil kali luar (outer product) di EMT hanya dapat dilakukan antar vektor. Hal ini otomatis karena
			bahasa matriksnya. Satu vektor harus berupa vektor kolom dan yang lainnya berupa vektor baris.
		\end{eulercomment}
		\begin{eulerprompt}
			>(1:5)*(1:5)'
		\end{eulerprompt}
		\begin{euleroutput}
			        1         2         3         4         5 
			        2         4         6         8        10 
			        3         6         9        12        15 
			        4         8        12        16        20 
			        5        10        15        20        25 
		\end{euleroutput}
		\begin{eulercomment}
			Dalam PDF pengantar R terdapat contoh yang menghitung distribusi dari $ab-cd$ untuk $a,b,c,d$ yang
			dipilih secara acak dari 0 hingga $n$. Solusi di R adalah membentuk matriks 4-dimensi dan menjalankan
			\texttt{table()} di atasnya.
			
			Tentu saja, hal ini dapat dicapai dengan perulangan. Namun perulangan tidak efektif di EMT maupun R.
			Di EMT, kita dapat menulis perulangan dalam bahasa C dan itu akan menjadi solusi tercepat.
			
			Namun kita ingin meniru perilaku R. Untuk ini, kita perlu meratakan (flatten) perkalian $ab$ dan
			membuat matriks dari $ab-cd$.
		\end{eulercomment}
		\begin{eulerprompt}
			>a=0:6; b=a'; p=flatten(a*b); q=flatten(p-p'); ...
			>u=sort(unique(q)); f=getmultiplicities(u,q); ...
			>statplot(u,f,"h"):
		\end{eulerprompt}
		\eulerimg{25}{images/Tugas 7-047.png}
		\begin{eulercomment}
			Selain multiplisitas eksak, EMT dapat menghitung frekuensi dalam vektor.
		\end{eulercomment}
		\begin{eulerprompt}
			>getfrequencies(q,-50:10:50)
		\end{eulerprompt}
		\begin{euleroutput}
			[0,  23,  132,  316,  602,  801,  333,  141,  53,  0]
		\end{euleroutput}
		\begin{eulercomment}
			Cara termudah untuk memplot ini sebagai distribusi adalah sebagai berikut.
		\end{eulercomment}
		\begin{eulerprompt}
			>plot2d(q,distribution=11):
		\end{eulerprompt}
		\eulerimg{25}{images/Tugas 7-048.png}
		\begin{eulercomment}
			Namun dimungkinkan juga untuk menghitung terlebih dahulu jumlah kemunculan dalam interval yang
			dipilih. Tentu saja, perintah berikut menggunakan \texttt{getfrequencies()} secara internal.
			
			Karena fungsi \texttt{histo()} mengembalikan frekuensi, kita perlu menskalakannya agar integral di bawah
			grafik batang bernilai 1.
		\end{eulercomment}
		\begin{eulerprompt}
			>{x,y}=histo(q,v=-55:10:55); y=y/sum(y)/differences(x); ...
			>plot2d(x,y,>bar,style="/"):
		\end{eulerprompt}
		\eulerimg{25}{images/Tugas 7-049.png}
		
		\eulerheading{Daftar (Lists)}
		\begin{eulercomment}
			EMT memiliki dua jenis daftar. Yang pertama adalah daftar global yang dapat diubah (mutable), dan
			yang lainnya adalah tipe daftar yang tidak dapat diubah (immutable). Kita tidak membahas daftar
			global di sini.
			
			Tipe daftar yang tidak dapat diubah disebut koleksi (collection) di EMT. Ini berperilaku seperti
			struktur (struct) di C, tetapi elemen-elemennya hanya diberi nomor dan tidak diberi nama.
		\end{eulercomment}
		\begin{eulerprompt}
			>L={{"Fred","Flintstone",40,[1990,1992]}}
		\end{eulerprompt}
		\begin{euleroutput}
			Fred
			Flintstone
			40
			[1990,  1992]
		\end{euleroutput}
		\begin{eulercomment}
			Saat ini elemen-elemen tersebut tidak memiliki nama, meskipun nama dapat diatur untuk tujuan khusus.
			Elemen-elemen tersebut diakses menggunakan nomor.
		\end{eulercomment}
		\begin{eulerprompt}
			>(L[4])[2]
		\end{eulerprompt}
		\begin{euleroutput}
			1992
		\end{euleroutput}
		
		\eulerheading{Masukan dan Keluaran File (Membaca dan Menulis Data)}
		\begin{eulercomment}
			Anda akan sering ingin mengimpor matriks data dari sumber lain ke EMT. Tutorial ini menjelaskan
			banyak cara untuk mencapainya. Fungsi sederhananya adalah \texttt{writematrix()} dan \texttt{readmatrix()}.
			
			Mari kita demonstrasikan cara membaca dan menulis vektor bilangan real ke dalam sebuah file.
		\end{eulercomment}
		\begin{eulerprompt}
			>a=random(1,100); mean(a), dev(a),
		\end{eulerprompt}
		\begin{euleroutput}
			0.49815
			0.28037
		\end{euleroutput}
		\begin{eulercomment}
			Untuk menulis data ke file, kita menggunakan fungsi \texttt{writematrix()}.
			
			Karena pengantar ini kemungkinan besar berada di direktori di mana pengguna tidak memiliki akses
			tulis, kita menulis data ke direktori home pengguna. Untuk notebook milik sendiri, hal ini tidak
			diperlukan, karena file data akan ditulis ke direktori yang sama.
		\end{eulercomment}
		\begin{eulerprompt}
			>filename="test.dat";
		\end{eulerprompt}
		\begin{eulercomment}
			Sekarang kita menulis vektor kolom \texttt{a'} ke dalam file. Ini menghasilkan satu angka di setiap baris
			file.
		\end{eulercomment}
		\begin{eulerprompt}
			>writematrix(a',filename);
		\end{eulerprompt}
		\begin{eulercomment}
			Untuk membaca data, kita menggunakan \texttt{readmatrix()}.
		\end{eulercomment}
		\begin{eulerprompt}
			>a=readmatrix(filename)';
		\end{eulerprompt}
		\begin{eulercomment}
			Dan hapus file tersebut.
		\end{eulercomment}
		\begin{eulerprompt}
			>fileremove(filename);
			>mean(a), dev(a),
		\end{eulerprompt}
		\begin{euleroutput}
			0.49815
			0.28037
		\end{euleroutput}
		\begin{eulercomment}
			Fungsi \texttt{writematrix()} atau \texttt{writetable()} dapat dikonfigurasi untuk bahasa lain.
			
			Misalnya, jika Anda menggunakan sistem Indonesia (tanda desimal dengan koma), Excel Anda memerlukan
			nilai dengan koma desimal yang dipisahkan oleh titik koma dalam file CSV (default-nya adalah nilai
			yang dipisahkan koma). File \texttt{"test.csv"} berikut akan muncul di folder Anda saat ini.
		\end{eulercomment}
		\begin{eulerprompt}
			>filename="test.csv"; ...
			>writematrix(random(5,3),file=filename,separator=",");
		\end{eulerprompt}
		\begin{eulercomment}
			Anda sekarang dapat membuka file ini secara langsung menggunakan Excel berbahasa Indonesia.
		\end{eulercomment}
		\begin{eulerprompt}
			>fileremove(filename);
		\end{eulerprompt}
		\begin{eulercomment}
			Terkadang kita memiliki string dengan token seperti berikut.
		\end{eulercomment}
		\begin{eulerprompt}
			>s1:="f m m f m m m f f f m m f";  ...
			>s2:="f f f m m f f";
		\end{eulerprompt}
		\begin{eulercomment}
			Untuk melakukan tokenisasi ini, kita mendefinisikan vektor token.
		\end{eulercomment}
		\begin{eulerprompt}
			>tok:=["f","m"]
		\end{eulerprompt}
		\begin{euleroutput}
			f
			m
		\end{euleroutput}
		\begin{eulercomment}
			Kemudian kita dapat menghitung berapa kali setiap token muncul dalam string, dan memasukkan hasilnya
			ke dalam tabel.
		\end{eulercomment}
		\begin{eulerprompt}
			>M:=getmultiplicities(tok,strtokens(s1))_ ...
			>  getmultiplicities(tok,strtokens(s2));
		\end{eulerprompt}
		\begin{eulercomment}
			Tulis tabel dengan judul kolom token.
		\end{eulercomment}
		\begin{eulerprompt}
			>writetable(M,labc=tok,labr=1:2,wc=8)
		\end{eulerprompt}
		\begin{euleroutput}
			               f       m
			       1       6       7
			       2       5       2
		\end{euleroutput}
		\begin{eulercomment}
			Untuk statistika, EMT dapat membaca dan menulis tabel.
		\end{eulercomment}
		\begin{eulerprompt}
			>file="test.dat"; open(file,"w"); ...
			>writeln("A,B,C"); writematrix(random(3,3)); ...
			>close();
		\end{eulerprompt}
		\begin{eulercomment}
			Isi file tersebut terlihat seperti ini.
		\end{eulercomment}
		\begin{eulerprompt}
			>printfile(file)
		\end{eulerprompt}
		\begin{euleroutput}
			A,B,C
			0.7003664386138074,0.1875530821001213,0.3262339279660414
			0.5926249243193858,0.1522927283984059,0.368140583062521
			0.8065535209872989,0.7265910840408142,0.7332619844597152
		\end{euleroutput}
		\begin{eulercomment}
			Fungsi \texttt{readtable()} dalam bentuk paling sederhananya dapat membaca ini dan mengembalikan koleksi nilai
			beserta baris judulnya.
		\end{eulercomment}
		\begin{eulerprompt}
			>L=readtable(file,>list);
		\end{eulerprompt}
		\begin{eulercomment}
			Koleksi ini dapat dicetak dengan \texttt{writetable()} ke notebook, atau ke sebuah file.
		\end{eulercomment}
		\begin{eulerprompt}
			>writetable(L,wc=10,dc=5)
		\end{eulerprompt}
		\begin{euleroutput}
			         A         B         C
			   0.70037   0.18755   0.32623
			   0.59262   0.15229   0.36814
			   0.80655   0.72659   0.73326
		\end{euleroutput}
		\begin{eulercomment}
			Matriks nilai adalah elemen pertama dari \texttt{L}. Perhatikan bahwa \texttt{mean()} di EMT menghitung nilai rata-rata
			dari baris-baris suatu matriks.
		\end{eulercomment}
		\begin{eulerprompt}
			>mean(L[1])
		\end{eulerprompt}
		\begin{euleroutput}
			  0.40472 
			  0.37102 
			  0.75547 
		\end{euleroutput}
		
		\eulerheading{File CSV}
		\begin{eulercomment}
			Pertama, mari kita tulis sebuah matriks ke dalam file. Untuk keluarannya, kita membuat file di
			direktori kerja saat ini.
		\end{eulercomment}
		\begin{eulerprompt}
			>file="test.csv";  ...
			>M=random(3,3); writematrix(M,file);
		\end{eulerprompt}
		\begin{eulercomment}
			Berikut adalah isi dari file ini.
		\end{eulercomment}
		\begin{eulerprompt}
			>printfile(file)
		\end{eulerprompt}
		\begin{euleroutput}
			0.8221197733097619,0.821531098722547,0.7771240608094004
			0.8482947121863489,0.3237767724883862,0.6501422353377985
			0.1482301827518109,0.3297459716109594,0.6261901074210923
		\end{euleroutput}
		\begin{eulercomment}
			File CSV ini dapat dibuka pada sistem berbahasa Inggris ke dalam Excel dengan klik ganda. Jika Anda
			mendapatkan file seperti ini pada sistem berbahasa Jerman, Anda perlu mengimpor data ke Excel dengan
			memperhatikan tanda titik desimal.
			
			Namun titik desimal juga merupakan format default untuk EMT. Anda dapat membaca matriks dari file
			dengan \texttt{readmatrix()}.
		\end{eulercomment}
		\begin{eulerprompt}
			>readmatrix(file)
		\end{eulerprompt}
		\begin{euleroutput}
			  0.82212   0.82153   0.77712 
			  0.84829   0.32378   0.65014 
			  0.14823   0.32975   0.62619 
		\end{euleroutput}
		\begin{eulercomment}
			Dimungkinkan untuk menulis beberapa matriks ke dalam satu file. Perintah \texttt{open()} dapat membuka file
			untuk ditulis dengan parameter \texttt{"w"}. Nilai default-nya adalah \texttt{"r"} untuk membaca.
		\end{eulercomment}
		\begin{eulerprompt}
			>open(file,"w"); writematrix(M); writematrix(M'); close();
		\end{eulerprompt}
		\begin{eulercomment}
			Matriks-matriks tersebut dipisahkan oleh baris kosong. Untuk membaca matriksnya, buka file tersebut
			dan panggil \texttt{readmatrix()} beberapa kali.
		\end{eulercomment}
		\begin{eulerprompt}
			>open(file); A=readmatrix(); B=readmatrix(); A==B, close();
		\end{eulerprompt}
		\begin{euleroutput}
			        1         0         0 
			        0         1         0 
			        0         0         1 
		\end{euleroutput}
		\begin{eulercomment}
			Di Excel atau program lembar kerja (spreadsheet) serupa, Anda dapat mengekspor matriks sebagai CSV
			(comma separated values). Di Excel 2007, gunakan ``save as'' dan ``other formats'', lalu pilih ``CSV''.
			Pastikan tabel saat ini hanya berisi data yang ingin Anda ekspor.
			
			Berikut adalah contohnya.
		\end{eulercomment}
		\begin{eulerprompt}
			>printfile("excel-data.csv")
		\end{eulerprompt}
		\begin{euleroutput}
			0;1000;1000
			1;1051,271096;1072,508181
			2;1105,170918;1150,273799
			3;1161,834243;1233,67806
			4;1221,402758;1323,129812
			5;1284,025417;1419,067549
			6;1349,858808;1521,961556
			7;1419,067549;1632,31622
			8;1491,824698;1750,6725
			9;1568,312185;1877,610579
			10;1648,721271;2013,752707
		\end{euleroutput}
		\begin{eulercomment}
			Seperti yang Anda lihat, sistem berbahasa Jerman saya menggunakan titik koma sebagai pemisah dan koma
			desimal. Anda dapat mengubahnya di pengaturan sistem atau di Excel, tetapi hal ini tidak diperlukan
			untuk membaca matriks ke dalam EMT.
			
			Cara termudah untuk membaca ini ke dalam Euler adalah \texttt{readmatrix()}. Semua koma digantikan oleh titik
			dengan parameter \texttt{>comma}. Untuk CSV berbahasa Inggris, cukup hilangkan parameter ini.
		\end{eulercomment}
		\begin{eulerprompt}
			>M=readmatrix("excel-data.csv",>comma)
		\end{eulerprompt}
		\begin{euleroutput}
			        0      1000      1000 
			        1    1051.3    1072.5 
			        2    1105.2    1150.3 
			        3    1161.8    1233.7 
			        4    1221.4    1323.1 
			        5      1284    1419.1 
			        6    1349.9      1522 
			        7    1419.1    1632.3 
			        8    1491.8    1750.7 
			        9    1568.3    1877.6 
			       10    1648.7    2013.8 
		\end{euleroutput}
		\begin{eulercomment}
			Mari kita plot ini.
		\end{eulercomment}
		\begin{eulerprompt}
			>plot2d(M'[1],M'[2:3],>points,color=[red,green]'):
		\end{eulerprompt}
		\eulerimg{25}{images/Tugas 7-050.png}
		\begin{eulercomment}
			Terdapat cara yang lebih mendasar untuk membaca data dari file. Anda dapat membuka file dan membaca
			angka-angkanya baris demi baris. Fungsi \texttt{getvectorline()} akan membaca angka dari satu baris data.
			Secara default, fungsi ini mengharapkan titik desimal. Namun fungsi ini juga dapat menggunakan koma
			desimal, jika Anda memanggil \texttt{setdecimaldot(",")} sebelum menggunakan fungsi ini.
			
			Fungsi berikut adalah contoh untuk ini. Fungsi akan berhenti di akhir file atau pada baris kosong.
		\end{eulercomment}
		\begin{eulerprompt}
			>function myload (file) ...
		\end{eulerprompt}
		\begin{eulerudf}
			open(file);
			M=[];
			repeat
			   until eof();
			   v=getvectorline(3);
			   if length(v)>0 then M=M_v; else break; endif;
			end;
			return M;
			close(file);
			endfunction
		\end{eulerudf}
		\begin{eulerprompt}
			>myload(file)
		\end{eulerprompt}
		\begin{euleroutput}
			  0.82212   0.82153   0.77712 
			  0.84829   0.32378   0.65014 
			  0.14823   0.32975   0.62619 
		\end{euleroutput}
		\begin{eulercomment}
			Dimungkinkan juga untuk membaca semua angka dalam file tersebut dengan \texttt{getvector()}.
		\end{eulercomment}
		\begin{eulerprompt}
			>open(file); v=getvector(10000); close(); redim(v[1:9],3,3)
		\end{eulerprompt}
		\begin{euleroutput}
			  0.82212   0.82153   0.77712 
			  0.84829   0.32378   0.65014 
			  0.14823   0.32975   0.62619 
		\end{euleroutput}
		\begin{eulercomment}
			Dengan demikian sangat mudah untuk menyimpan vektor nilai, satu nilai di setiap baris, dan membaca
			kembali vektor ini.
		\end{eulercomment}
		\begin{eulerprompt}
			>v=random(1000); mean(v)
		\end{eulerprompt}
		\begin{euleroutput}
			0.50303
		\end{euleroutput}
		\begin{eulerprompt}
			>writematrix(v',file); mean(readmatrix(file)')
		\end{eulerprompt}
		\begin{euleroutput}
			0.50303
		\end{euleroutput}
		
		\eulerheading{Menggunakan Tabel}
		\begin{eulercomment}
			Tabel dapat digunakan untuk membaca atau menulis data numerik. Sebagai contoh, kita menulis tabel
			dengan judul baris dan kolom ke sebuah file.
		\end{eulercomment}
		\begin{eulerprompt}
			>file="test.tab"; M=random(3,3);  ...
			>open(file,"w");  ...
			>writetable(M,separator=",",labc=["one","two","three"]);  ...
			>close(); ...
			>printfile(file)
		\end{eulerprompt}
		\begin{euleroutput}
			one,two,three
			      0.09,      0.39,      0.86
			      0.39,      0.86,      0.71
			       0.2,      0.02,      0.83
		\end{euleroutput}
		\begin{eulercomment}
			Ini dapat diimpor ke dalam Excel.
			
			Untuk membaca file di EMT, kita menggunakan \texttt{readtable()}.
		\end{eulercomment}
		\begin{eulerprompt}
			>{M,headings}=readtable(file,>clabs); ...
			>writetable(M,labc=headings)
		\end{eulerprompt}
		\begin{euleroutput}
			       one       two     three
			      0.09      0.39      0.86
			      0.39      0.86      0.71
			       0.2      0.02      0.83
		\end{euleroutput}
		
		\eulerheading{Menganalisis Baris}
		\begin{eulercomment}
			Anda bahkan dapat mengevaluasi setiap baris secara manual. Misalkan kita memiliki baris dengan format
			berikut.
		\end{eulercomment}
		\begin{eulerprompt}
			>line="2020-11-03,Tue,1'114.05"
		\end{eulerprompt}
		\begin{euleroutput}
			2020-11-03,Tue,1'114.05
		\end{euleroutput}
		\begin{eulercomment}
			Pertama kita dapat melakukan tokenisasi pada baris tersebut.
		\end{eulercomment}
		\begin{eulerprompt}
			>vt=strtokens(line)
		\end{eulerprompt}
		\begin{euleroutput}
			2020-11-03
			Tue
			1'114.05
		\end{euleroutput}
		\begin{eulercomment}
			Kemudian kita dapat mengevaluasi setiap elemen baris menggunakan evaluasi yang sesuai.
		\end{eulercomment}
		\begin{eulerprompt}
			>day(vt[1]),  ...
			>indexof(["mon","tue","wed","thu","fri","sat","sun"],tolower(vt[2])),  ...
			>strrepl(vt[3],"'","")()
		\end{eulerprompt}
		\begin{euleroutput}
			7.3816e+05
			2
			1114
		\end{euleroutput}
		\begin{eulercomment}
			Dengan menggunakan ekspresi reguler (regular expressions), dimungkinkan untuk mengekstrak hampir
			semua informasi dari sebaris data.
			
			Misalkan kita memiliki baris berikut dalam sebuah dokumen HTML.
		\end{eulercomment}
		\begin{eulerprompt}
			>line="<tr><td>1145.45</td><td>5.6</td><td>-4.5</td><tr>"
		\end{eulerprompt}
		\begin{euleroutput}
			<tr><td>1145.45</td><td>5.6</td><td>-4.5</td><tr>
		\end{euleroutput}
		\begin{eulercomment}
			Untuk mengekstrak ini, kita menggunakan ekspresi reguler yang mencari:
			\begin{itemize}
				\item kurung tutup \texttt{>},
				\item string apa pun yang tidak berisi tanda kurung dengan sub-kecocokan \texttt{"(...)"},
				\item kurung buka dan kurung tutup menggunakan solusi terpendek,
				\item lagi-lagi string apa pun yang tidak berisi tanda kurung,
				\item dan kurung buka \texttt{<}.
			\end{itemize}
			
			Ekspresi reguler agak sulit dipelajari tetapi sangat ampuh.
		\end{eulercomment}
		\begin{eulerprompt}
			>{pos,s,vt}=strxfind(line,">([^<>]+)<.+?>([^<>]+)<");
		\end{eulerprompt}
		\begin{eulercomment}
			Hasilnya adalah posisi kecocokan, string yang cocok, dan vektor string untuk sub-kecocokan.
		\end{eulercomment}
		\begin{eulerprompt}
			>for k=1:length(vt); vt[k](), end;
		\end{eulerprompt}
		\begin{euleroutput}
			1145.5
			5.6
		\end{euleroutput}
		\begin{eulercomment}
			Berikut adalah fungsi yang membaca semua item numerik di antara \texttt{<td>} dan \texttt{</td>}.
		\end{eulercomment}
		\begin{eulerprompt}
			>function readtd (line) ...
		\end{eulerprompt}
		\begin{eulerudf}
			v=[]; cp=0;
			repeat
			   {pos,s,vt}=strxfind(line,"<td.*?>(.+?)</td>",cp);
			   until pos==0;
			   if length(vt)>0 then v=v|vt[1]; endif;
			   cp=pos+strlen(s);
			end;
			return v;
			endfunction
		\end{eulerudf}
		\begin{eulerprompt}
			>readtd(line+"<td>non-numerical</td>")
		\end{eulerprompt}
		\begin{euleroutput}
			1145.45
			5.6
			-4.5
			non-numerical
		\end{euleroutput}
		
		\eulerheading{Membaca dari Web}
		\begin{eulercomment}
			Situs web atau file dengan URL dapat dibuka di EMT dan dapat dibaca baris demi baris.
			
			Pada contoh ini, kita membaca versi saat ini dari situs EMT. Kita menggunakan ekspresi reguler untuk
			memindai \texttt{"Version ..."} dalam sebuah judul.
		\end{eulercomment}
		\begin{eulerprompt}
			>function readversion () ...
		\end{eulerprompt}
		\begin{eulerudf}
			urlopen("http://www.euler-math-toolbox.de/Programs/Changes.html");
			repeat
			  until urleof();
			  s=urlgetline();
			  k=strfind(s,"Version ",1);
			  if k>0 then substring(s,k,strfind(s,"<",k)-1), break; endif;
			end;
			urlclose();
			endfunction
		\end{eulerudf}
		\begin{eulerprompt}
			>readversion
		\end{eulerprompt}
		\begin{euleroutput}
			Version 2024-01-12
		\end{euleroutput}
		
		\eulerheading{Masukan dan Keluaran Variabel}
		\begin{eulercomment}
			Anda dapat menulis variabel dalam bentuk definisi Euler ke sebuah file atau ke baris perintah.
		\end{eulercomment}
		\begin{eulerprompt}
			>writevar(pi,"mypi");
		\end{eulerprompt}
		\begin{euleroutput}
			mypi = 3.141592653589793;
		\end{euleroutput}
		\begin{eulercomment}
			Sebagai pengujian, kita membuat file Euler di direktori kerja EMT.
		\end{eulercomment}
		\begin{eulerprompt}
			>file="test.e"; ...
			>writevar(random(2,2),"M",file); ...
			>printfile(file,3)
		\end{eulerprompt}
		\begin{euleroutput}
			M = [ ..
			0.5991820585590205, 0.7960280262224293;
			0.5167243983231363, 0.2996684599070898];
		\end{euleroutput}
		\begin{eulercomment}
			Sekarang kita dapat memuat file tersebut. File ini akan mendefinisikan matriks \texttt{M}.
		\end{eulercomment}
		\begin{eulerprompt}
			>load(file); show M,
		\end{eulerprompt}
		\begin{euleroutput}
			M = 
			  0.59918   0.79603 
			  0.51672   0.29967 
		\end{euleroutput}
		\begin{eulercomment}
			Omong-omong, jika \texttt{writevar()} digunakan pada suatu variabel, fungsi tersebut akan mencetak definisi
			variabel beserta nama variabel ini.
		\end{eulercomment}
		\begin{eulerprompt}
			>writevar(M); writevar(inch$)
		\end{eulerprompt}
		\begin{euleroutput}
			M = [ ..
			0.5991820585590205, 0.7960280262224293;
			0.5167243983231363, 0.2996684599070898];
			inch$ = 0.0254;
		\end{euleroutput}
		\begin{eulercomment}
			Kita juga dapat membuka file baru atau menambahkan (append) ke file yang sudah ada. Pada contoh ini
			kita menambahkan ke file yang dibuat sebelumnya.
		\end{eulercomment}
		\begin{eulerprompt}
			>open(file,"a"); ...
			>writevar(random(2,2),"M1"); ...
			>writevar(random(3,1),"M2"); ...
			>close();
			>load(file); show M1; show M2;
		\end{eulerprompt}
		\begin{euleroutput}
			M1 = 
			  0.30287   0.15372 
			   0.7504   0.75401 
			M2 = 
			  0.27213 
			 0.053211 
			  0.70249 
		\end{euleroutput}
		\begin{eulercomment}
			Untuk menghapus file apa pun gunakan \texttt{fileremove()}.
		\end{eulercomment}
		\begin{eulerprompt}
			>fileremove(file);
		\end{eulerprompt}
		\begin{eulercomment}
			Vektor baris dalam sebuah file tidak memerlukan koma, jika setiap angka berada pada baris baru. Mari
			kita buat file seperti itu, menulis setiap baris satu per satu dengan \texttt{writeln()}.
		\end{eulercomment}
		\begin{eulerprompt}
			>open(file,"w"); writeln("M = ["); ...
			>for i=1 to 5; writeln(""+random()); end; ...
			>writeln("];"); close(); ...
			>printfile(file)
		\end{eulerprompt}
		\begin{euleroutput}
			M = [
			0.344851384551
			0.0807510017715
			0.876519562911
			0.754157709472
			0.688392638934
			];
		\end{euleroutput}
		\begin{eulerprompt}
			>load(file); M
		\end{eulerprompt}
		\begin{euleroutput}
			[0.34485,  0.080751,  0.87652,  0.75416,  0.68839]
		\end{euleroutput}
		
		\eulerheading{Materi Lainnya}
		\begin{eulercomment}
			Topik lanjutan yang dibahas pada bagian-bagian berikut meliputi:
			\begin{itemize}
				\item Ukuran statistika deskriptif tambahan (jangkauan, modus, dan metode pergeseran nilai taksiran),
				\item Uji kesesuaian distribusi normal dengan kertas probabilitas normal,
				\item Interval kepercayaan untuk rata-rata dan varians,
				\item Uji independensi korelasi dan interval kepercayaan koefisien regresi,
				\item Regresi linear multidimensi (berganda) dan uji ekuivalensi,
				\item Algoritma pembangkit bilangan acak dan transformasi invers,
				\item Aplikasi metode Monte Carlo untuk integral lipat dan persamaan diferensial parsial (masalah nilai batas Laplace).
			\end{itemize}
		\end{eulercomment}
		
		\eulerheading{Ukuran Statistika Deskriptif Tambahan}
		\begin{eulercomment}
			Selain rata-rata, median, dan simpangan baku, ukuran deskriptif lain yang sering digunakan adalah
			jangkauan (range) dan modus. Jangkauan $R$ didefinisikan sebagai selisih antara nilai maksimum dan
			minimum:
		\end{eulercomment}
		\begin{eulerformula}
			\[
			R = x_{\max} - x_{\min}.
			\]
		\end{eulerformula}
		\begin{eulercomment}
			Sedangkan modus adalah nilai data yang muncul dengan frekuensi paling besar. Berikut adalah contoh
			pengukuran induksi celah udara pada 15 pengeras suara (dalam satuan Tesla).
		\end{eulercomment}
		\begin{eulerprompt}
			>x = [1.01, 1.02, 1.00, 0.98, 0.99, 1.00, 0.99, 1.01, ...
			>     1.01, 1.00, 1.00, 1.00, 1.02, 1.00, 1.01];
		\end{eulerprompt}
		\begin{eulercomment}
			Kita dapat menghitung rata-rata, varians sampel $s^2$, median, serta jangkauan $R$ dari sampel tersebut.
		\end{eulercomment}
		\begin{eulerprompt}
			>xbar = mean(x), s2 = dev(x)^2, med = median(x), R = max(x) - min(x),
		\end{eulerprompt}
		\begin{euleroutput}
			1.0027
			0.00012095
			1
			0.04
		\end{euleroutput}
		\begin{eulercomment}
			Untuk mencari modus, kita menghitung multiplisitas setiap nilai unik dalam sampel menggunakan fungsi
			\texttt{unique()} dan \texttt{getmultiplicities()}.
		\end{eulercomment}
		\begin{eulerprompt}
			>u = unique(x); freq = getmultiplicities(u, x);
			>modus = u[nonzeros(freq == max(freq))]
		\end{eulerprompt}
		\begin{euleroutput}
			1
		\end{euleroutput}
		\begin{eulercomment}
			Untuk bilangan dengan banyak digit yang hanya berbeda pada beberapa digit terakhir, perhitungan
			rata-rata dan varians dapat disederhanakan menggunakan nilai taksiran awal $x_0$ melalui transformasi
		\end{eulercomment}
		\begin{eulerformula}
			\[
			z_i = x_i - x_0
			\]
		\end{eulerformula}
		\begin{eulerformula}
			\[
			\bar{x} = x_0 + \bar{z}, \qquad s^2 = \frac{\sum_{i=1}^n z_i^2 - n(\bar{x}-x_0)^2}{n-1}.
			\]
		\end{eulerformula}
		\begin{eulercomment}
			Mari kita terapkan metode ini dengan memilih
		\end{eulercomment}
		\begin{eulerformula}
			\[
			x_0 = 1.00.
			\]
		\end{eulerformula}
		\begin{eulerprompt}
			>x0 = 1.00; z = x - x0; n = length(x);
			>xbarkode = x0 + mean(z)
		\end{eulerprompt}
		\begin{euleroutput}
			1.0027
		\end{euleroutput}
		\begin{eulerprompt}
			>s2kode = (sum(z^2) - n*(xbarkode - x0)^2) / (n - 1)
		\end{eulerprompt}
		\begin{euleroutput}
			0.00012095
		\end{euleroutput}
		
		\eulerheading{Uji Kesesuaian dengan Kertas Probabilitas Normal}
		\begin{eulercomment}
			Untuk memeriksa secara grafis apakah data sampel berasal dari distribusi normal, kita dapat
			menggunakan prinsip kertas probabilitas normal (normal probability paper). Jika $X \sim N(\mu, \sigma^2)$,
			maka fungsi distribusi kumulatifnya memenuhi hubungan linear:
		\end{eulercomment}
		\begin{eulerformula}
			\[
			z = \Phi^{-1}(F(x)) = \frac{x - \mu}{\sigma}.
			\]
		\end{eulerformula}
		\begin{eulercomment}
			Dengan demikian, jika frekuensi kumulatif relatif $F_j$ ditransformasikan melalui invers distribusi
			normal standar $\Phi^{-1}(F_j)$ dan diplot terhadap batas atas kelas $\xi_j$,
			titik-titik tersebut akan terletak mendekati sebuah garis lurus.
			
			Berikut adalah data $n = 125$ pengukuran yang dikelompokkan ke dalam 11 kelas dari 50 hingga 270 dengan
			lebar kelas $h = 20$.
		\end{eulercomment}
		\begin{eulerprompt}
			>xi = 70:20:270; // batas atas setiap kelas
			>h = [1, 1, 2, 9, 15, 22, 30, 27, 9, 6, 3]; // frekuensi absolut
			>F = cumsum(h) / sum(h); // frekuensi kumulatif relatif
		\end{eulerprompt}
		\begin{eulercomment}
			Karena $\Phi^{-1}(1) = \infty$, kita memplot 10 batas atas kelas pertama terhadap nilai $z_j = \Phi^{-1}(F_j)$
			menggunakan fungsi \texttt{invnormaldis()}.
		\end{eulercomment}
		\begin{eulerprompt}
			>z = invnormaldis(F[1:10], 0, 1);
			>p = polyfit(xi[1:10], z, 1); // pencocokan garis lurus z = p[1] + p[2]*x
			>plot2d(xi[1:10], z, >points, style="x", color=red, ...
			>  xl="Batas Atas Kelas", yl="z = invnormaldis(F)"); ...
			>plot2d("evalpoly(x,p)", 60, 260, color=blue, >add):
		\end{eulerprompt}
		\eulerimg{25}{images/Tugas 7-061.png}
		\begin{eulercomment}
			Dari kemiringan dan titik potong garis lurus
		\end{eulercomment}
		\begin{eulerformula}
			\[
			z = -\frac{\mu}{\sigma} + \frac{1}{\sigma}x,
			\]
		\end{eulerformula}
		\begin{eulercomment}
			kita dapat menaksir parameter $\mu$ dan $\sigma$ secara langsung.
		\end{eulercomment}
		\begin{eulerprompt}
			>sigmaest = 1 / p[2], muest = -p[1] / p[2],
		\end{eulerprompt}
		\begin{euleroutput}
			38.986
			174.88
		\end{euleroutput}
		
		\eulerheading{Interval Kepercayaan untuk Rata-Rata dan Varians}
		\eulersubheading{1. Interval Kepercayaan Rata-Rata (Varians Populasi $\sigma^2$ Diketahui)}
		\begin{eulercomment}
			Jika sampel acak berukuran $n$ diambil dari populasi dengan varians $\sigma^2$ yang diketahui dan tingkat
			signifikansi $\alpha$, batas kepercayaan untuk nilai harapan $\mu$ adalah:
		\end{eulercomment}
		\begin{eulerformula}
			\[
			\bar{x} \pm \frac{\sigma}{\sqrt{n}} z_{1-\alpha/2}.
			\]
		\end{eulerformula}
		\begin{eulercomment}
			Berikut adalah fungsi untuk menghitung interval kepercayaan tersebut.
		\end{eulercomment}
		\begin{eulerprompt}
			>function cimeanknown (x, sigma, alpha=0.05) ...
		\end{eulerprompt}
		\begin{eulerudf}
			n = length(x);
			zcrit = invnormaldis(1 - alpha/2, 0, 1);
			err = zcrit * sigma / sqrt(n);
			return [mean(x) - err, mean(x) + err];
			endfunction
		\end{eulerudf}
		\begin{eulerprompt}
			>cimeanknown([0.842, 0.846, 0.835, 0.839, 0.843, 0.838], 0.004, 0.05)
		\end{eulerprompt}
		\begin{euleroutput}
			[0.8373,  0.8437]
		\end{euleroutput}
		
		\eulersubheading{2. Interval Kepercayaan Rata-Rata (Varians Populasi $\sigma^2$ Tidak Diketahui)}
		\begin{eulercomment}
			Jika populasi berdistribusi normal dan $\sigma^2$ tidak diketahui, kita menggantinya dengan simpangan
			baku sampel $s$ dan menggunakan kuantil distribusi-$t$ dengan derajat kebebasan $m = n - 1$:
		\end{eulercomment}
		\begin{eulerformula}
			\[
			\bar{x} \pm \frac{s}{\sqrt{n}} t_{1-\alpha/2; \, n-1}.
			\]
		\end{eulerformula}
		\begin{eulercomment}
			Mari kita terapkan pada sampel 6 hasil pengukuran dengan tingkat signifikansi $\alpha = 5\%$ dan $\alpha = 1\%$.
		\end{eulercomment}
		\begin{eulerprompt}
			>function cimeanunknown (x, alpha=0.05) ...
		\end{eulerprompt}
		\begin{eulerudf}
			n = length(x);
			tcrit = invtdis(1 - alpha/2, n - 1);
			err = tcrit * dev(x) / sqrt(n);
			return [mean(x) - err, mean(x) + err];
			endfunction
		\end{eulerudf}
		\begin{eulerprompt}
			>sampel = [0.842, 0.846, 0.835, 0.839, 0.843, 0.838];
			>mean(sampel), dev(sampel),
		\end{eulerprompt}
		\begin{euleroutput}
			0.8405
			0.003937
		\end{euleroutput}
		\begin{eulerprompt}
			>cimeanunknown(sampel, 0.05) // tingkat kepercayaan 95%
		\end{eulerprompt}
		\begin{euleroutput}
			[0.83637,  0.84463]
		\end{euleroutput}
		\begin{eulerprompt}
			>cimeanunknown(sampel, 0.01) // tingkat kepercayaan 99%
		\end{eulerprompt}
		\begin{euleroutput}
			[0.83402,  0.84698]
		\end{euleroutput}
		
		\eulersubheading{3. Interval Kepercayaan untuk Varians $\sigma^2$}
		\begin{eulercomment}
			Untuk populasi berdistribusi normal, besaran $(n-1)s^2/\sigma^2$ mengikuti distribusi $\chi^2$ dengan derajat kebebasan $m = n - 1$. Interval kepercayaan $(1-\alpha)$ untuk
			$\sigma^2$ diberikan oleh:
		\end{eulercomment}
		\begin{eulerformula}
			\[
			\frac{(n-1)s^2}{\chi^2_{1-\alpha/2; \, n-1}} \le \sigma^2 \le \frac{(n-1)s^2}{\chi^2_{\alpha/2; \, n-1}}.
			\]
		\end{eulerformula}
		\begin{eulerprompt}
			>function civar (x, alpha=0.05) ...
		\end{eulerprompt}
		\begin{eulerudf}
			n = length(x);
			s2 = dev(x)^2;
			chilower = invchidis(alpha/2, n - 1);
			chiupper = invchidis(1 - alpha/2, n - 1);
			return [(n - 1)*s2 / chiupper, (n - 1)*s2 / chilower];
			endfunction
		\end{eulerudf}
		\begin{eulerprompt}
			>civar(sampel, 0.05) // interval kepercayaan 95% untuk varians
		\end{eulerprompt}
		\begin{euleroutput}
			[0,  9.3237e-05]
		\end{euleroutput}
		\begin{eulerprompt}
			>sqrt(civar(sampel, 0.05)) // interval kepercayaan 95% untuk simpangan baku
		\end{eulerprompt}
		\begin{euleroutput}
			[0.0024575,  0.009656]
		\end{euleroutput}
		
		\eulerheading{Uji Independensi Korelasi dan Interval Koefisien Regresi}
		\eulersubheading{1. Uji Independensi Dua Variabel ($H_0: \rho_{XY} = 0$)}
		\begin{eulercomment}
			Untuk menguji apakah dua variabel $X$ dan $Y$ yang berdistribusi normal bivariat saling independen
			($\rho_{XY} = 0$), kita menghitung statistik uji-$t$ dari koefisien korelasi empiris $r_{xy}$:
		\end{eulercomment}
		\begin{eulerformula}
			\[
			t = \frac{r_{xy}\sqrt{n-2}}{\sqrt{1 - r_{xy}^2}},
			\]
		\end{eulerformula}
		\begin{eulercomment}
			yang berdistribusi-$t$ dengan derajat kebebasan $m = n - 2$.
		\end{eulercomment}
		\begin{eulerprompt}
			>function correltest (x, y) ...
		\end{eulerprompt}
		\begin{eulerudf}
			n = length(x);
			r = correl(x, y);
			tstat = r * sqrt(n - 2) / sqrt(1 - r^2);
			pval = 2 * (1 - tdis(abs(tstat), n - 2));
			return {r, tstat, pval};
			endfunction
		\end{eulerudf}
		\begin{eulerprompt}
			>x = 1:10; y = [2, 3, 1, 5, 6, 3, 7, 8, 9, 8];
			>{r, tstat, pval} = correltest(x, y);
			>r, tstat, pval,
		\end{eulerprompt}
		\begin{euleroutput}
			0.87175
			5.0325
			0.001011
		\end{euleroutput}
		\begin{eulercomment}
			Karena nilai-$p$ jauh lebih kecil dari 0.05, hipotesis nol bahwa $X$ dan $Y$ tidak berkorelasi ditolak.
		\end{eulercomment}
		
		\eulersubheading{2. Interval Kepercayaan untuk Koefisien Regresi Linear}
		\begin{eulercomment}
			Pada garis regresi
		\end{eulercomment}
		\begin{eulerformula}
			\[
			y(x) = \tilde{a} + \tilde{b}x,
			\]
		\end{eulerformula}
		\begin{eulercomment}
			batas kepercayaan $(1-\alpha)$ untuk kemiringan $\tilde{b}$ dan titik potong $\tilde{a}$ masing-masing
			dihitung dengan
		\end{eulercomment}
		\begin{eulerformula}
			\[
			\tilde{b} \pm t_{1-\alpha/2; \, n-2} \frac{s_y \sqrt{1 - r_{xy}^2}}{s_x \sqrt{n-2}}, \qquad \tilde{a} \pm t_{1-\alpha/2; \, n-2} \frac{s_y \sqrt{1 - r_{xy}^2} \sqrt{\sum x_i^2}}{s_x \sqrt{n-2} \sqrt{n}}.
			\]
		\end{eulerformula}
		\begin{eulerprompt}
			>function cireg (x, y, alpha=0.05) ...
		\end{eulerprompt}
		\begin{eulerudf}
			n = length(x);
			p = polyfit(x, y, 1); // p[1] = a, p[2] = b
			r = correl(x, y);
			sx = dev(x); sy = dev(y);
			tcrit = invtdis(1 - alpha/2, n - 2);
			sefactor = sy * sqrt(1 - r^2) / (sx * sqrt(n - 2));
			errb = tcrit * sefactor;
			erra = tcrit * sefactor * sqrt(sum(x^2) / n);
			return [p[1] - erra, p[1] + erra; p[2] - errb, p[2] + errb];
			endfunction
		\end{eulerudf}
		\begin{eulerprompt}
			>cireg(x, y, 0.05) // baris 1: interval untuk a, baris 2: interval untuk b
		\end{eulerprompt}
		\begin{euleroutput}
			  -1.5757    3.0424 
			  0.43999    1.1843 
		\end{euleroutput}
		
		\eulerheading{Regresi Multidimensi (Berganda)}
		\begin{eulercomment}
			Jika variabel tak bebas $Y$ bergantung secara linear pada beberapa fungsi basis $g_j(\underline{x})$,
			model regresinya ditulis sebagai
		\end{eulercomment}
		\begin{eulerformula}
			\[
			y = \sum_{j=0}^s a_j g_j(\underline{x}).
			\]
		\end{eulerformula}
		\begin{eulercomment}
			Dengan metode kuadrat terkecil, vektor taksiran koefisien $\tilde{\underline{a}}$
			diperoleh dengan menyelesaikan sistem persamaan normal:
		\end{eulercomment}
		\begin{eulerformula}
			\[
			G^T G \tilde{\underline{a}} = G^T \underline{f}.
			\]
		\end{eulerformula}
		\begin{eulercomment}
			Mari kita selesaikan contoh model
		\end{eulercomment}
		\begin{eulerformula}
			\[
			\tilde{f}(x_1, x_2) = a_0 + a_1 x_1 + a_2 x_2
			\]
		\end{eulerformula}
		\begin{eulercomment}
			untuk 4 titik pengamatan.
		\end{eulercomment}
		\begin{eulerprompt}
			>x1 = [5, 3, 5, 3]';
			>x2 = [0.5, 0.5, 0.3, 0.3]';
			>f = [1.5, 3.5, 6.2, 3.2]';
			>G = ones(4, 1) | x1 | x2
		\end{eulerprompt}
		\begin{euleroutput}
			        1         5       0.5 
			        1         3       0.5 
			        1         5       0.3 
			        1         3       0.3 
		\end{euleroutput}
		\begin{eulercomment}
			Kita menyusun matriks persamaan normal $G^T G$ dan vektor sisi kanan $G^T \underline{f}$,
			lalu menyelesaikannya menggunakan operator \texttt{\textbackslash}.
		\end{eulercomment}
		\begin{eulerprompt}
			>ahat = (G'.G) \ (G'.f)
		\end{eulerprompt}
		\begin{euleroutput}
			        7 
			     0.25 
			      -11 
		\end{euleroutput}
		\begin{eulercomment}
			Hasilnya adalah
		\end{eulercomment}
		\begin{eulerformula}
			\[
			\tilde{a}_0 = 7.0, \quad \tilde{a}_1 = 0.25, \quad \text{dan} \quad \tilde{a}_2 = -11.0.
			\]
		\end{eulerformula}
		\begin{eulercomment}
			Kita juga dapat menghitung jumlah kuadrat sisaan (residual sum of squares) serta melakukan uji
			ekuivalensi dengan distribusi-$F$ untuk membandingkan model penuh dengan model tereduksi.
		\end{eulercomment}
		\begin{eulerprompt}
			>resfull = f - G.ahat;
			>SSRfull = sum((resfull^2)')
		\end{eulerprompt}
		\begin{euleroutput}
			6.25
		\end{euleroutput}
		\begin{eulercomment}
			Misalkan kita ingin menguji apakah koefisien $a_2$ dapat diabaikan ($H_0: a_2 = 0$), sehingga model
			tereduksi hanya menggunakan $x_1$.
		\end{eulercomment}
		\begin{eulerprompt}
			>Gred = ones(4, 1) | x1;
			>ared = (Gred'.Gred) \ (Gred'.f);
			>resred=f-Gred.ared;
			>SSRred = sum(((resred)^2)');
			>df1 = cols(G) - cols(Gred); // selisih jumlah parameter = 1
			>df2 = rows(G) - cols(G); // derajat kebebasan sisaan = 4 - 3 = 1
			>Fstat = ((SSRred - SSRfull) / df1) / (SSRfull / df2)
		\end{eulerprompt}
		\begin{euleroutput}
			0.7744
		\end{euleroutput}
		\begin{eulerprompt}
			>pvalF = 1 - fdis(Fstat, df1, df2)
		\end{eulerprompt}
		\begin{euleroutput}
			0.54058
		\end{euleroutput}
		\begin{eulercomment}
			Karena nilai-$p$ lebih besar dari 0.05, penghapusan variabel $x_2$ tidak memperburuk kecocokan model
			secara signifikan pada sampel kecil ini.
		\end{eulercomment}
		
		\eulerheading{Algoritma Pembangkit Bilangan Acak}
		\eulersubheading{1. Metode Kuadrat Tengah (J. von Neumann)}
		\begin{eulercomment}
			Metode ini dimulai dari bilangan desimal $z_0 \in (0, 1)$
			dengan $2n$ digit, mengkuadratkannya sehingga menghasilkan $4n$ digit, lalu mengambil $2n$ digit di bagian
			tengah sebagai bilangan acak berikutnya.
		\end{eulercomment}
		\begin{eulerprompt}
			>function midsquare (z0, k, ndigits=4) ...
		\end{eulerprompt}
		\begin{eulerudf}
			res = zeros(1, k);
			z = z0;
			half = ndigits / 2;
			for i = 1 to k;
			  res[i] = z;
			  val = floor(z^2 * 10^(ndigits + half));
			  z = mod(val, 10^ndigits) / 10^ndigits;
			end;
			return res;
			endfunction
		\end{eulerudf}
		\begin{eulerprompt}
			>midsquare(0.1234, 6, 4)
		\end{eulerprompt}
		\begin{euleroutput}
			[0.1234,  0.5227,  0.3215,  0.3362,  0.303,  0.1809]
		\end{euleroutput}
		\begin{eulercomment}
			Metode ini jarang digunakan dalam praktik modern karena deretnya dapat mengalami degenerasi menuju
			nol.
		\end{eulercomment}
		
		\eulersubheading{2. Metode Kongruensi Linear (Linear Congruential Generator)}
		\begin{eulercomment}
			Metode yang jauh lebih luas digunakan adalah rekursi kongruensi:
		\end{eulercomment}
		\begin{eulerformula}
			\[
			z_{i+1} \equiv c \cdot z_i \pmod{m}, \qquad u_i = \frac{z_i}{m}.
			\]
		\end{eulerformula}
		\begin{eulerprompt}
			>function lcg (z0, c, m, k) ...
		\end{eulerprompt}
		\begin{eulerudf}
			u = zeros(1, k);
			z = z0;
			for i = 1 to k;
			  z = mod(c * z, m);
			  u[i] = z / m;
			end;
			return u;
			endfunction
		\end{eulerudf}
		\begin{eulerprompt}
			>lcg(123, 16807, 2^31 - 1, 8)
		\end{eulerprompt}
		\begin{euleroutput}
			[0.00096264,  0.17915,  0.93945,  0.41397,  0.53037,  0.93198,  0.78649,  0.50036]
		\end{euleroutput}
		
		\eulersubheading{3. Metode Transformasi Invers ($F^{-1}(\xi)$)}
		\begin{eulercomment}
			Jika $\xi \sim U(0, 1)$ adalah bilangan acak berdistribusi seragam dan $F(x)$ adalah fungsi distribusi kumulatif kontinu yang
			memiliki invers, maka
		\end{eulercomment}
		\begin{eulerformula}
			\[
			\eta = F^{-1}(\xi) \quad \text{mengikuti distribusi } F(x).
			\]
		\end{eulerformula}
		\begin{eulercomment}
			Sebagai contoh, untuk distribusi eksponensial
		\end{eulercomment}
		\begin{eulerformula}
			\[
			F(x) = 1 - e^{-\lambda x},
			\]
		\end{eulerformula}
		\begin{eulercomment}
			fungsi inversnya adalah
		\end{eulercomment}
		\begin{eulerformula}
			\[
			F^{-1}(\xi) = -\frac{1}{\lambda}\ln(1 - \xi).
			\]
		\end{eulerformula}
		\begin{eulerprompt}
			>lambda = 2; xi = random(1, 5000);
			>eta = -log(1 - xi) / lambda;
			>plot2d(eta, distribution=25, style="\/"); ...
			>plot2d("2*exp(-2*x)", 0, max(eta), color=red, thickness=2, >add):
		\end{eulerprompt}
		\eulerimg{25}{images/Tugas 7-088.png}
		
		\eulerheading{Aplikasi Metode Monte Carlo dalam Matematika Numerik}
		\eulersubheading{1. Evaluasi Integral Tunggal}
		\begin{eulercomment}
			Untuk menghitung integral tentu
		\end{eulercomment}
		\begin{eulerformula}
			\[
			I = \int_0^1 g(x)\,dx
			\]
		\end{eulerformula}
		\begin{eulercomment}
			dengan $0 \le g(x) \le 1$, terdapat dua pendekatan Monte Carlo:
			\begin{itemize}
				\item \textbf{Metode Frekuensi Relatif (Hit-or-Miss):} Membangkitkan $n$ pasang titik acak $(\xi_i, \eta_i) \in [0, 1]^2$ dan menghitung proporsi titik $m/n$ yang memenuhi $\eta_i \le g(\xi_i)$.
				\item \textbf{Metode Nilai Rata-Rata:} Menghitung rata-rata fungsi pada titik acak seragam $\xi_i \in [0, 1]$:
			\end{itemize}
		\end{eulercomment}
		\begin{eulerformula}
			\[
			\int_0^1 g(x)\,dx \approx \frac{1}{n}\sum_{i=1}^n g(\xi_i).
			\]
		\end{eulerformula}
		\begin{eulercomment}
			Mari kita bandingkan kedua metode tersebut untuk menghitung
		\end{eulercomment}
		\begin{eulerformula}
			\[
			\int_0^1 \frac{4}{1 + x^2}\,dx = \pi \approx 3.14159265.
			\]
		\end{eulerformula}
		\begin{eulercomment}
			Karena nilai maksimum integran adalah 4, kita menskalakan fungsi dengan $g(x) = \frac{1}{1+x^2} \in [0, 1]$
			lalu mengalikan hasilnya dengan 4.
		\end{eulercomment}
		\begin{eulerprompt}
			>n = 100000;
			>xi = random(1, n); eta = random(1, n);
			>g = 1 / (1 + xi^2);
			>pihitormiss = 4 * sum(eta <= g) / n
		\end{eulerprompt}
		\begin{euleroutput}
			3.1475
		\end{euleroutput}
		\begin{eulerprompt}
			>pimeanvalue = 4 * mean(g)
		\end{eulerprompt}
		\begin{euleroutput}
			3.1429
		\end{euleroutput}
		\begin{eulercomment}
			Metode nilai rata-rata umumnya memberikan galat yang lebih kecil dibandingkan metode frekuensi
			relatif untuk jumlah sampel $n$ yang sama.
		\end{eulercomment}
		
		\eulersubheading{2. Evaluasi Integral Lipat Dua}
		\begin{eulercomment}
			Untuk menghitung volume
		\end{eulercomment}
		\begin{eulerformula}
			\[
			V = \iint_S h(x, y)\,dx\,dy
			\]
		\end{eulerformula}
		\begin{eulercomment}
			pada domain $S = [a, b] \times [c, d]$ dengan $0 \le h(x, y) \le e$, 
			kita membangkitkan tripel bilangan acak seragam $(u_i, v_i, w_i) \in [0, 1]^3$
			dan menghitung proporsi titik di bawah permukaan $h(x, y)$.
			
			Sebagai contoh, kita hitung
		\end{eulercomment}
		\begin{eulerformula}
			\[
			V = \int_0^1 \int_0^1 (x^2 + y^2)\,dx\,dy = \frac{2}{3}
			\]
		\end{eulerformula}
		\begin{eulercomment}
			(dengan batas atas $e = 2$).
		\end{eulercomment}
		\begin{eulerprompt}
			>n = 200000;
			>u = random(1, n); v = random(1, n); w = random(1, n);
			>e = 2;
			>Vhit = e * sum(w <= (u^2 + v^2) / e) / n
		\end{eulerprompt}
		\begin{euleroutput}
			0.66031
		\end{euleroutput}
		\begin{eulerprompt}
			>Vmean = mean(u^2 + v^2)
		\end{eulerprompt}
		\begin{euleroutput}
			0.66526
		\end{euleroutput}
		
		\eulersubheading{3. Penyelesaian Masalah Nilai Batas Laplace dengan Random Walk}
		\begin{eulercomment}
			Metode Monte Carlo dapat digunakan untuk menghampiri solusi persamaan diferensial parsial Laplace
		\end{eulercomment}
		\begin{eulerformula}
			\[
			\Delta u = \frac{\partial^2 u}{\partial x^2} + \frac{\partial^2 u}{\partial y^2} = 0
			\]
		\end{eulerformula}
		\begin{eulercomment}
			pada domain $G$ dengan syarat batas $u(x, y) = f(x, y)$ pada $\Gamma$.
			
			Domain $G$ didiskretisasi menggunakan kisi persegi dengan langkah $h = 1$. Sebuah partikel dimulai dari
			titik interior $P(x_0, y_0)$ dan bergerak secara acak ke salah satu dari empat titik tetangganya
			masing-masing dengan probabilitas $1/4$ hingga menyentuh titik batas $R_i \in \Gamma$. Nilai $u(x_0, y_0)$
			dihampiri oleh rata-rata nilai batas $f(R_i)$ tempat jalan acak berhenti:
		\end{eulercomment}
		\begin{eulerformula}
			\[
			v(x_0, y_0) \approx \frac{1}{n} \sum_{k=1}^n f(R^{(k)}).
			\]
		\end{eulerformula}
		\begin{eulercomment}
			Kita uji metode ini pada kisi persegi $0 \le x \le 4, \; 0 \le y \le 4$
			dengan fungsi batas harmonik $f(x, y) = x + 2y$. Solusi eksak pada titik interior $(2, 1)$ adalah $u(2, 1) = 2 + 2(1) = 4$.
		\end{eulercomment}
		\begin{eulerprompt}
			>function laplacerw (x0, y0, nwalks)
		\end{eulerprompt}
		\begin{eulerudf}
			total = 0;
			for k = 1 to nwalks;
			  x = x0; y = y0;
			  repeat
			    until (x == 0 || x == 4 || y == 0 || y == 4);
			    dir = intrandom(1, 1, 4)[1];
			    if dir == 1 then x = x + 1;
			    elseif dir == 2 then x = x - 1;
			    elseif dir == 3 then y = y + 1;
			    else y = y - 1;
			    endif;
			  end;
			  total = total + (x + 2*y); // evaluasi syarat batas f(x, y) = x + 2y
			end;
			return total / nwalks;
			endfunction
		\end{eulerudf}
		\begin{eulerprompt}
			>laplacerw(2, 1, 5000)
		\end{eulerprompt}
		\begin{euleroutput}
			4.0116
		\end{euleroutput}
	\end{eulernotebook}
\end{document}